\documentclass{report}
\usepackage{bm}
\newcommand{\xvec}{\bm{x}}
\usepackage{pgfplots}                        %
\pgfplotsset{samples=50}

\newcommand{\yvec}{\mathbf{y}}

%
\usepackage[
  a4paper,
  margin=0.7in,
  includehead,
  headheight=14pt,
  headsep=8pt
]{geometry}   

\usepackage{setspace}                        %
\usepackage{fancyhdr}                        %
\usepackage{titlesec}                        %
\titleformat{\chapter}[hang]
  {\normalfont\Huge\bfseries}
  {\thechapter}{1em}{}

\titlespacing*{\chapter}
  {0pt}{-20pt}{20pt}
\usepackage{longtable}
\usepackage[T1]{fontenc}
\usepackage[utf8]{inputenc}
\usepackage{lmodern}                         %
\usepackage{algorithmicx}                    %
\usepackage{algorithm}                      %
\usepackage{algpseudocode}                  %
\usepackage{float}                          %
\usepackage{microtype}                       %
\usepackage{amsmath,amssymb,amsthm}
\usepackage{graphicx}
\usepackage{booktabs}
\usepackage{multirow}
\usepackage{hyperref}
\usepackage{listings}

%
\usepackage{tikz}
\usetikzlibrary{arrows.meta, positioning, calc}
\usepackage[acronym]{glossaries}
\usepackage{amsfonts}
\usepackage{enumitem}
\usepackage[numbers,sort&compress]{natbib}

%
\usepackage{xcolor}

\definecolor{codebg}{HTML}{F6F8FA}

%
\theoremstyle{plain}
\newtheorem{theorem}{Theorem}[chapter]
\newtheorem{proposition}[theorem]{Proposition}
\newtheorem{lemma}[theorem]{Lemma}
\newtheorem{corollary}[theorem]{Corollary}
\newtheorem{claim}[theorem]{Claim}

%
\theoremstyle{remark}
\newtheorem{remark}[theorem]{Remark}

%
\onehalfspacing                               %
\setlength{\parindent}{1.2em}                 %
\setlength{\parskip}{0.2em}                   %
\numberwithin{equation}{section}              %

%
\titleformat{\section}{\normalfont\Large\bfseries}{\thesection}{0.8em}{}
\titleformat{\subsection}{\normalfont\large\bfseries}{\thesubsection}{0.6em}{}
\titleformat{\subsubsection}{\normalfont\normalsize\bfseries}{\thesubsubsection}{0.6em}{}
\titlespacing*{\section}{0pt}{1.1ex plus 0.4ex minus 0.2ex}{0.8ex}
\titlespacing*{\subsection}{0pt}{1.0ex plus 0.4ex minus 0.2ex}{0.6ex}
\titlespacing*{\subsubsection}{0pt}{0.9ex plus 0.3ex minus 0.2ex}{0.5ex}

%
\setcounter{tocdepth}{1}
\setcounter{secnumdepth}{3}

%
\hypersetup{
  colorlinks=true,
  linkcolor=black,
  citecolor=black,
  urlcolor=blue
}

%
\pagestyle{fancy}
\fancyhf{}                                   %
\lhead{\leftmark}                            %
\rhead{\thepage}                             %
\renewcommand{\headrulewidth}{0.4pt}
\renewcommand{\footrulewidth}{0pt}

%
%
%
\makeatletter
\newenvironment{breakablealgorithm}
  {%
    \begin{center}
      \refstepcounter{algorithm}%
      \hrule height .8pt depth 0pt \kern 2pt%
      \renewcommand{\caption}[2][\relax]{%
        {\raggedright\textbf{Algorithm~\thealgorithm.} ##2\par}%
        \ifx\relax##1\relax
        \else
          \addcontentsline{loa}{algorithm}{\protect\numberline{\thealgorithm}##1}%
        \fi
        \kern 2pt\hrule\kern 2pt
      }%
  }
  {%
      \kern 2pt\hrule\relax%
    \end{center}
  }
\makeatother

%
\algrenewcommand\algorithmicindent{1.0em}

%
\newcommand{\M}{\mathcal{M}}            %
\newcommand{\gradM}{\nabla_{\!\M}}      %
\newcommand{\divM}{\nabla_{\!\M}\!\cdot}%
\newcommand{\LapM}{\Delta_{\!\M}}       %
\newcommand{\bvec}{\mathbf{b}}
\newcommand{\Gmat}{\mathbf{G}}
\newcommand{\Wmat}{\mathbf{W}}
\newcommand{\uvec}{\mathbf{u}}
\newcommand{\pvec}{\mathbf{p}}
\newcommand{\feedbackgraphic}[2]{%
  \IfFileExists{#2}{%
    \includegraphics[width=#1]{#2}%
  }{%
    \fbox{%
      \parbox[c][0.16\textheight][c]{#1}{%
        \centering Figure omitted in single-file draft%
      }%
    }%
  }%
}

%
\newacronym{wsindy}{WSINDy}{Weak-form Sparse Identification of Nonlinear Dynamics}
\newacronym{mvar}{MVAR}{Multivariate Autoregression}
\newacronym{mstls}{MSTLS}{Multi-Step Thresholded Least Squares}
\newacronym{efrom}{EF-ROM}{Equation-Free Reduced-Order Model}
\newacronym{lstms}{LSTMs}{Long Short-Term Memory}
\newacronym{pod}{POD}{Proper Orthogonal Decomposition}
\newacronym{svd}{SVD}{Singular Value Decomposition}

\makeglossaries

%
\title{Interpretable Weak-Form PDE Discovery for Interacting-Agent Densities\\
from Noisy Trajectories: Benchmarking WSINDy Against Equation-Free Reduced-Order Models}

\author{Eugenio Macias}
\date{} %

\def\feedbackdraft{1}

%
%
%
\begin{document}

%
%
\begin{titlepage}
\centering

\vspace*{1.2in}

{\LARGE\bfseries
Interpretable Weak-Form PDE Discovery for Interacting-Agent Densities\\
from Noisy Trajectories: Benchmarking WSINDy Against Equation-Free Reduced-Order Models
\par}

\vspace{0.8in}

{\Large Eugenio Macias\par}

\vspace{0.6in}

{\large
Honors Thesis\\
Applied Mathematics (APMA)\\
Brown University
\par}

\vspace{0.5in}

{\large Advisor: Prof.\ Bj\"orn Sandstede\par}

\vfill

%
%

\end{titlepage}
 
%
\pagenumbering{roman}

%
\begin{abstract}
This thesis studies interpretable weak-form PDE discovery for macroscopic density evolution inferred from noisy microscopic trajectory data generated by discrete Vicsek alignment with Morse attraction--repulsion forces, and benchmarks it against equation-free reduced-order models. On the benchmark side, density snapshots are shift-aligned via FFT-based cross-correlation on the periodic torus, transformed and compressed by \gls{pod}, evolved in latent space by two contrasting time-steppers, a linear multivariate autoregression (\gls{mvar}) and a nonlinear long short-term memory (LSTM) network, and then lifted back to physical density fields. This contrast tests whether the aligned low-dimensional dynamics are effectively linear after reduction or whether materially nonlinear sequence modelling is needed. On the interpretable side, a weak-form sparse identification framework (\gls{wsindy}) discovers coupled PDEs for the density~$\rho$ and polarization flux fields~$p_x$, $p_y$ from a physically motivated candidate library, with sparsity enforced by modified sequential thresholded least squares (MSTLS) and discovered models integrated forward via an exponential time-differencing scheme (ETDRK4). Performance is evaluated across seven dynamical regimes, four initial-condition families, and multiple noise levels. The central question is whether WSINDy can recover compact, interpretable flux laws whose forecasting accuracy is competitive with the equation-free reduced-order benchmarks. All code, configuration files, and analysis scripts are publicly available at \url{https://github.com/eugenio-macias/wsindy-manifold}.
\end{abstract}
 
\clearpage
\tableofcontents

%
\chapter*{Notation and Symbols}
\addcontentsline{toc}{chapter}{Notation and Symbols}

\noindent
Mathematical symbols and the corresponding Python variable names used
throughout the repository \texttt{wsindy-manifold}.

\setlength{\LTpre}{6pt}
\setlength{\LTpost}{0pt}
\renewcommand{\arraystretch}{1.15}

\newcommand{\cv}[1]{{\ttfamily\footnotesize #1}}   %
\begin{longtable}{%
  p{0.20\textwidth}
  p{0.48\textwidth}
  p{0.24\textwidth}}
\toprule
\textbf{Symbol} & \textbf{Name / Role} & \textrm{\textbf{Code variable}} \\
\midrule
\endfirsthead

\toprule
\textbf{Symbol} & \textbf{Name / Role} & \textrm{\textbf{Code variable}} \\
\midrule
\endhead

\midrule
\multicolumn{3}{r}{\emph{(Continued on next page)}}\\
\bottomrule
\endfoot

\bottomrule
\endlastfoot

%
\multicolumn{3}{l}{\textbf{Geometry and Domains}}\\
$\Omega=[0,L_x]\times[0,L_y]$ & Rectangular computational domain & \cv{Lx, Ly} \\
$\Gamma$ & Boundary of $\Omega$ & \cv{bc} \\
$n_x, n_y$ & Density-grid resolution & \cv{DENSITY\_NX, DENSITY\_NY} \\
\addlinespace

%
\multicolumn{3}{l}{\textbf{Simulation Parameters}}\\
$N$ & Number of particles & \cv{N} \\
$T_{\mathrm{sim}}$ & Simulation duration (s) & \cv{T} \\
$\Delta t$ & Simulation time-step & \cv{dt} \\
$v_0$ & Particle speed & \cv{speed} \\
$\eta$ & Angular-noise strength & \cv{eta} \\
$R$ & Alignment interaction radius & \cv{R} \\
\addlinespace

%
\multicolumn{3}{l}{\textbf{Fields and Densities}}\\
$\rho(t,x,y)$ & Agent density field & \cv{rho} \\
$p_x(t,x,y),\;p_y(t,x,y)$ & Polarisation flux (velocity-weighted KDE) & \cv{px, py} \\
$\Phi(t,x,y)$ & Morse potential field ($W*\rho$) & \cv{Phi} \\
$W(r)$ & Morse interaction kernel & \cv{W} \\
$C_a,\;C_r$ & Attraction / repulsion strengths & \cv{Ca, Cr} \\
$\ell_a,\;\ell_r$ & Attraction / repulsion length scales & \cv{la, lr} \\
\addlinespace

%
\multicolumn{3}{l}{\textbf{Differential Operators}}\\
$\partial_t$ & Time derivative & --- \\
$\nabla,\;\nabla\!\cdot$ & Spatial gradient / divergence & --- \\
$\Delta = \partial_{xx}+\partial_{yy}$ & Laplacian & \cv{lap} \\
$\Delta^2$ & Biharmonic operator & \cv{bilap} \\
\addlinespace

%
\multicolumn{3}{l}{\textbf{Shift Alignment}}\\
$\rho_{\mathrm{ref}}$ & Reference density for alignment & \cv{ref} \\
$(\Delta y_t,\Delta x_t)$ & Per-frame integer pixel shifts & \cv{shifts} \\
$\tilde{\rho}(t,x,y)$ & Aligned density field & \cv{aligned} \\
$\hat{\rho}_{\mathrm{ref}},\;\hat{\rho}_{\mathrm{target}}$ & FFTs of reference / target density & \cv{F\_ref, F\_target} \\
\addlinespace

%
\multicolumn{3}{l}{\textbf{Data and Density Estimation}}\\
$\{\mathbf{x}_i(t)\}_{i=1}^N$ & Agent trajectories & \cv{traj} \\
$\{\mathbf{v}_i(t)\}_{i=1}^N$ & Agent velocities & \cv{vel} \\
$\sigma_h$ & KDE bandwidth (Gaussian kernel) & \cv{bandwidth} \\
$g(\cdot)$ & Density transform (raw / log / sqrt / meansub) & \cv{density\_transform} \\
$\epsilon$ & Numerical offset for log/sqrt & \cv{density\_transform\_eps} \\
\addlinespace

%
\multicolumn{3}{l}{\textbf{POD / SVD}}\\
$\mathbf{X}$ & Snapshot matrix (training densities, flattened) & \cv{X\_all} \\
$\overline{\mathbf{X}}$ & Temporal mean field & \cv{X\_mean} \\
$\mathbf{X}'$ & Mean-centred snapshots & \cv{X\_centered} \\
$\mathbf{U}\boldsymbol{\Sigma}\mathbf{V}^{\!\top}$ & SVD factorisation of $\mathbf{X}'$ & \cv{U, S, Vt} \\
$\mathbf{U}_r$ & POD basis (first $r$ columns of $\mathbf{U}$) & \cv{U\_r} \\
$r$ & Truncation rank (number of retained modes) & \cv{R\_POD} \\
$\xi$ & Energy threshold (default 0.995) & \cv{TARGET\_ENERGY} \\
$\mathbf{y}(t)$ & Latent coordinates ($\mathbf{U}_r^{\!\top}\mathbf{X}'$) & \cv{X\_latent} \\
$M$ & Number of training runs & \cv{M} \\
$\Delta_r$ & Temporal subsampling factor & \cv{ROM\_SUBSAMPLE} \\
\addlinespace

%
\multicolumn{3}{l}{\textbf{MVAR Forecasting}}\\
$p$ & VAR lag order & \cv{P\_LAG} \\
$\alpha$ & Ridge regularisation parameter & \cv{RIDGE\_ALPHA} \\
$\mathbf{A}^{(1)},\dots,\mathbf{A}^{(p)}$ & VAR coefficient matrices & \cv{A\_matrices} \\
$\mathbf{c}$ & VAR intercept / bias & \cv{intercept} \\
$\mathbf{C}$ & Companion matrix ($pd\times pd$) & \cv{C} \\
$\rho_{\max}$ & Stability threshold (optional, usually off) & \cv{eigenvalue\_threshold} \\
$R^2_{\mathrm{train}}$ & Training coefficient of determination & \cv{r2\_train} \\
\addlinespace

%
\multicolumn{3}{l}{\textbf{WSINDy (Weak-form PDE Discovery)}}\\
$\psi(t,x,y)$ & Compactly supported test function & \cv{psi} \\
$\ell_t,\;\ell_x,\;\ell_y$ & Test-function half-widths (grid cells) & \cv{ellt, ellx, elly} \\
$p_t,\;p_x,\;p_y$ & Test-function polynomial exponents & \cv{pt, px, py} \\
$s_t,\;s_x,\;s_y$ & Query-point strides & \cv{stride\_t, stride\_x, stride\_y} \\
$\mathbf{G}$ & Weak-form design matrix & \cv{G} \\
$\mathbf{b}$ & Weak-form right-hand side ($\int\!\psi_t\,U$) & \cv{b} \\
$\mathbf{w}$ & Sparse coefficient vector & \cv{w} \\
$\hat{\lambda}$ & MSTLS sparsification threshold & \cv{lambda\_hat} \\
$\mathbf{s}$ & Column-norm scaling factors & \cv{col\_scale} \\
\textsc{MSTLS} & Multi-step thresholded least squares & \cv{mstls()} \\
\addlinespace

%
\multicolumn{3}{l}{\textbf{WSINDy Library}}\\
$(\mathcal{L}_j,\;f_j)$ & (Operator, feature) pair in library & \cv{library\_terms} \\
$f_j(U)$ & Nonlinear feature ($1, u, u^2, u^3, \dots$) & \cv{features} \\
$\mathcal{L}_j$ & Differential operator ($I$, $\partial_x$, $\partial_y$, $\Delta$, \dots) & \cv{operators} \\
$p_{\max}$ & Maximum polynomial power & \cv{max\_poly} \\
$\mathbf{m}$ & Binary active-term mask & \cv{active\_mask} \\
\addlinespace

%
\multicolumn{3}{l}{\textbf{ETDRK4 Time Integration}}\\
$\hat{L}(k_x,k_y)$ & Fourier multiplier --- linear part & \cv{L\_hat} \\
$\hat{N}(u)$ & Fourier multiplier --- nonlinear part & \cv{N\_hat} \\
$k_x,\;k_y$ & Wavenumber grids & \cv{KX, KY} \\
$E=e^{\Delta t\,\hat{L}}$ & Full-step exponential & \cv{E} \\
$E_2=e^{\Delta t\,\hat{L}/2}$ & Half-step exponential & \cv{E2} \\
$Q$ & Intermediate-stage $\varphi_1$ weight & \cv{Q} \\
$f_1,\;f_2,\;f_3$ & Cox--Matthews final-stage weights & \cv{f1, f2, f3} \\
\addlinespace

%
\multicolumn{3}{l}{\textbf{Metrics and Diagnostics}}\\
$R^2$ & Coefficient of determination & \cv{r2\_pod} \\
$R^2(t)$ & Time-resolved $R^2$ & \cv{r2\_time} \\
$\mathrm{RMSE}$ & Root-mean-squared error & \cv{rmse} \\
$\sigma_{\mathrm{rel}}$ & Relative error & \cv{rel\_error} \\
$\Delta M$ & Mass conservation error & \cv{mass\_error} \\
$\mathcal{T}_{\mathrm{train}},\;\mathcal{T}_{\mathrm{test}}$ & Training / testing time intervals & --- \\
\end{longtable}
 
\clearpage
\pagenumbering{arabic}

%
%
\chapter{Introduction}

We study whether an interpretable macroscopic PDE model can compete with black-box reduced-order surrogates when both are learned from noisy interacting-agent trajectories. The thesis benchmarks WSINDy against two POD-based latent models with deliberately different inductive biases: a linear multivariate autoregression (MVAR) and a nonlinear long short-term memory (LSTM) network. Starting from trajectories $\{x_i(t)\}_{i=1}^N \subset \Omega \subset \mathbb{R}^2$ generated by alignment and attraction--repulsion interactions, we coarse-grain to a smooth density field $\rho(x,t)$ and seek a predictive, physically interpretable evolution law. At the macroscopic level, density evolution obeys the conservation law $\partial_t \rho + \nabla\cdot \mathbf{j} = 0$, where $\mathbf{j}(x,t)$ is the agent flux. The modeling task is to identify a parameterized flux $\mathbf{j}_{\boldsymbol{\theta}}[\rho]$ from candidate mechanisms so that
\[
\partial_t \rho(x,t) + \nabla\cdot \mathbf{j}_{\boldsymbol{\theta}}[\rho](x,t) = 0
\]
holds on the observed data. In practice, we work in weak form to avoid unstable pointwise differentiation of noisy density fields.

Three challenges motivate the methodology. First, pointwise differentiation of noisy density fields yields unstable coefficient estimates. Second, the macroscopic flux depends not only on $\rho$ but also on local polarization, nonlocal interactions, and coefficients that vary with noise level, density regime, and coarse-graining scale. Third, coherent flock translation dominates the density movies, inflating the effective rank of the snapshot matrix under standard POD. A rich literature on symmetry reduction and shifted POD addresses this by factoring transport out before basis computation \cite{reiss2018shifted_pod,black2020projection,krah2024robust_spod,budanur2015slices,siminos2015periodic,sonday2011alignment}. We adopt a computationally efficient variant: each frame is registered to a common reference via FFT-based cross-correlation on the periodic torus, yielding integer-pixel shifts that remove translational drift without interpolation. The resulting aligned snapshot matrix admits a sharply faster singular-value decay, enabling faithful POD reconstruction with far fewer modes.

To address these challenges, we develop two complementary pipelines. Both start from the same upstream data generation: the same family of particle simulations, the same KDE-based macroscopic fields, and the same regime-specific train/test split. The first is an equation-free reduced-order model (EF-ROM) that compresses shift-aligned density snapshots via POD, evolves the latent coordinates with a learned time-stepper---either linear multivariate autoregression (MVAR) or a long short-term memory (LSTM) network---and lifts back to the physical field. Comparing MVAR and LSTM reveals whether the latent dynamics induced by Vicsek \cite{vicsek1995vicsek} alignment and D'Orsogna \cite{dorsogna2006morse} attraction--repulsion remain effectively linear after reduction.

The second pipeline targets interpretability via weak-form sparse identification (WSINDy). Here the learning object differs from the EF-ROM branch: WSINDy does not operate on the POD latent dataset, but instead remains in physical field space and identifies sparse continuous-time equations directly from weak-form systems built on the raw density field and auxiliary multi-fields. We formulate a coupled multi-field system: density $\rho$, polarization fluxes $p_x, p_y$ (from velocity-weighted KDE), and a Morse interaction potential $\Phi = W\ast\rho$ (via FFT convolution). The WSINDy library includes physically motivated candidate terms---continuity, diffusion, Morse forcing, nonlinear diffusion, alignment self-saturation, and density-driven pressure---yielding a three-equation discovery problem for $(\rho_t,\, p_{x,t},\, p_{y,t})$. Compactly supported polynomial test functions and integration by parts replace pointwise derivatives; sparsity is enforced via modified sequential thresholded least squares (MSTLS) with a dominant-balance criterion and automatic test-function width selection. Discovered PDEs are integrated forward using ETDRK4, which treats stiff linear terms exactly in Fourier space. The central question is whether WSINDy can recover compact, interpretable flux laws competitive with EF-ROM across noise levels, interaction strengths, and initial-condition families. Performance is evaluated via RMSE, forecast $R^2$, relative errors, forecast horizons, mass drift, and runtime; if the recovered PDEs are not yet competitive with the black-box benchmarks, the discrepancy is interpreted as evidence that the candidate library is incomplete, and residual structure is used to guide library augmentation.

\noindent Although all numerical results in this thesis are established on rectangular domains, the methodology is designed as groundwork for extension to nonplanar Riemannian manifolds $(\mathcal{M}, g)$. On such a surface, the flat-space operators used throughout must be replaced by their metric-aware counterparts: the Euclidean gradient $\nabla$ becomes the surface gradient $\nabla_{\!\mathcal{M}}$, the divergence $\nabla\cdot$ becomes $\nabla_{\!\mathcal{M}}\!\cdot$, and the Laplacian $\Delta$ becomes the Laplace--Beltrami operator $\Delta_{\!\mathcal{M}}$. Inner products in the weak form acquire the volume element $\sqrt{\det g}\,\mathrm{d}x$, convolutions such as $\Phi = W\ast\rho$ must respect geodesic distance rather than Euclidean distance, and the FFT-based shift alignment assumes periodicity on a flat torus that would need to be generalised to intrinsic registration on~$\mathcal{M}$. Once these substitutions are made, the weak-form library construction, MSTLS sparsity enforcement, and ETDRK4 integration all carry over in principle, providing a roadmap for geometry-consistent density fields PDE discovery beyond the plane.

\paragraph{Software.}The complete implementation accompanying this thesis---including the agent simulator, KDE coarse-graining, shift-aligned POD pipeline, MVAR and LSTM latent surrogates, WSINDy discovery engine, ETDRK4 integrator, and all YAML experiment configurations---is released as an open-source Python package at \url{https://github.com/eugenio-macias/wsindy-manifold}. Every figure and table in the following chapters can be reproduced from the scripts and configuration files therein.

\paragraph{Relation to prior work.}The EF-ROM component follows the multiscale restriction--lifting philosophy developed by Alvarez et al.~\cite{alvarez2025eqfree}. The alignment preprocessing draws on a rich literature on symmetry reduction, rotational equivariances, and shifted POD for transport-dominated problems \cite{sonday2011alignment,budanur2015slices,reiss2018shifted_pod,black2020projection,siminos2015periodic,krah2024robust_spod}; our implementation adapts the simplest variant---single-reference FFT cross-correlation on a periodic domain---which suffices for the dominant single-cluster translations observed in Vicsek dynamics. The PDE discovery component builds on the weak-form sparse identification methodology of Messenger and Bortz~\cite{messenger2021wsindy}. The simulator design and focus on systematic exploration of parameter space reflect current research directions in collective dynamics, where macroscopic laws are expected to remain stable across interaction regimes, noise levels, and densities. Finally, the framework is informed by growing interest in geometry-aware modeling, in which macroscopic laws must be consistent with the underlying spatial geometry on which collective motion occurs.

\begin{figure}[H]
\centering
\definecolor{sharedNavy}{HTML}{2C3E50}
\definecolor{efromTeal}{HTML}{D5F5E3}
\definecolor{efromLabel}{HTML}{117A65}
\definecolor{wsOrange}{HTML}{FDEBD0}
\definecolor{wsLabel}{HTML}{CA6F1E}
\begin{tikzpicture}[
    >=Stealth,
    node distance=0.5cm,
    every node/.style={font=\small},
    shared/.style={
      draw=none, rounded corners=4pt, minimum width=10cm, minimum height=0.8cm,
      fill=sharedNavy, text=white, font=\small\bfseries,
      align=center
    },
    efrom/.style={
      draw=none, rounded corners=4pt, minimum width=5.6cm, minimum height=0.8cm,
      fill=efromTeal, text=black, font=\small, align=center
    },
    wsindy/.style={
      draw=none, rounded corners=4pt, minimum width=5.6cm, minimum height=0.8cm,
      fill=wsOrange, text=black, font=\small, align=center
    },
    branchlabel/.style={font=\normalsize\bfseries\sffamily},
    desc/.style={font=\scriptsize\itshape, text=black!60},
    arr/.style={-Stealth, thick, draw=black!50},
    splitarr/.style={-Stealth, thick, draw=black!50}
  ]

  %
  \node[shared] (agents) {Agent Trajectories};
  \node[desc, below=0.03cm of agents] (agents-d) {Vicsek alignment $+$ Morse attraction--repulsion};

  \node[shared, below=0.4cm of agents-d] (kde) {KDE Density Fields};
  \node[desc, below=0.03cm of kde] (kde-d) {$\rho(x,t)$ on periodic grid via Gaussian kernel};

  \draw[arr] (agents) -- (kde);

  %
  \coordinate[below=0.65cm of kde-d] (split);

  %
  \node[branchlabel, anchor=south] at ([xshift=-3.1cm]split)
    (efrom-label) {\color{efromLabel}EF-ROM};
  \node[desc, below=0.0cm of efrom-label] (efrom-sub) {Part~I: Data-driven latent model};

  \node[efrom, below=0.3cm of efrom-sub] (align) {\textbf{Shift-Alignment}};
  \node[desc, below=0.03cm of align] (align-d) {FFT cross-correlation on periodic torus};

  \node[efrom, below=0.3cm of align-d] (pod) {\textbf{POD Compression}};
  \node[desc, below=0.03cm of pod] (pod-d) {$\sqrt{\cdot}\,$-transform, truncate to $r$ modes};

  %
  \coordinate[below=0.65cm of pod-d] (stepcentre);
  \node[efrom, minimum width=2.5cm] at ([xshift=-1.5cm]stepcentre) (mvar) {\textbf{MVAR}};
  \node[efrom, minimum width=2.5cm] at ([xshift=1.5cm]stepcentre)  (lstm) {\textbf{LSTM}};
  \node[desc, below=0.03cm of stepcentre, anchor=north] (step-d)
    {Latent time-stepping};

  \node[efrom, below=0.3cm of step-d] (lift) {\textbf{Lift to Density}};
  \node[desc, below=0.03cm of lift] (lift-d) {Inverse POD $+$ undo shift};

  %
  \draw[splitarr] (split) -| (align.north);
  \draw[arr] (align) -- (pod);
  \draw[arr] (pod.south) -- ++(0,-0.55cm) -| (mvar.north);
  \draw[arr] (pod.south) -- ++(0,-0.55cm) -| (lstm.north);
  \draw[arr] (mvar.south) |- ([yshift=0.15cm]lift.west) -- (lift.west);
  \draw[arr] (lstm.south) |- ([yshift=0.15cm]lift.east) -- (lift.east);

  %
  \node[branchlabel, anchor=south] at ([xshift=3.1cm]split)
    (ws-label) {\color{wsLabel}WSINDy};
  \node[desc, below=0.0cm of ws-label] (ws-sub) {Part~II: PDE discovery};

  \node[wsindy, below=0.3cm of ws-sub] (fields) {\textbf{Multi-field Extraction}};
  \node[desc, below=0.03cm of fields] (fields-d) {$\rho,\; p_x,\; p_y,\; \Phi = W\ast\rho$};

  \node[wsindy, below=0.3cm of fields-d] (lib) {\textbf{Weak-form Library}};
  \node[desc, below=0.03cm of lib] (lib-d) {Test functions $+$ candidate PDE terms};

  \node[wsindy, below=0.3cm of lib-d] (sparse) {\textbf{Sparse Regression}};
  \node[desc, below=0.03cm of sparse] (sparse-d) {MSTLS with dominant-balance criterion};

  \node[wsindy, below=0.3cm of sparse-d] (etdrk) {\textbf{ETDRK4 Integration}};
  \node[desc, below=0.03cm of etdrk] (etdrk-d) {Fourier-space time-stepping};

  %
  \draw[splitarr] (split) -| (fields.north);
  \draw[arr] (fields) -- (lib);
  \draw[arr] (lib) -- (sparse);
  \draw[arr] (sparse) -- (etdrk);

  %
  %
  \coordinate (merge) at ($(lift-d.south)!0.5!(etdrk-d.south) + (0,-0.7cm)$);
  \node[shared] at (merge) (eval) {Evaluation};
  \node[desc, below=0.03cm of eval] (eval-d)
    {RMSE, $R^2$, forecast horizon~$\tau$, mass drift, runtime};

  \draw[splitarr] (lift-d.south) |- ([yshift=0.35cm]eval.west) -- (eval.west);
  \draw[splitarr] (etdrk-d.south) |- ([yshift=0.35cm]eval.east) -- (eval.east);

\end{tikzpicture}
\caption{Overview of the two complementary modeling pipelines developed in this
  thesis.  Both branches share the same microscopic-to-macroscopic data
  generation (top) and are evaluated on common metrics (bottom).
  Left: the equation-free reduced-order model compresses aligned density
  snapshots into a low-dimensional latent space and evolves them with learned
  time-steppers.
  Right: WSINDy discovers interpretable PDE flux laws from multi-field data via
  weak-form sparse regression.}
\label{fig:pipeline_overview}
\end{figure}
 
%
\part{Equation-Free Reduced Order Modeling from Agent Simulations}

%
\chapter{Related Work and Background}
\label{ch:background}

\section{Collective motion models}

We adopt the simulation and data-generation methodology developed for crowd dynamics in \cite{alvarez2025eqfree} as the reference framework for the reduced-order modeling pipeline.
The discrete Vicsek simulator is implemented in \texttt{vicsek\_discrete.py}; force
evaluation in \texttt{morse.py}; boundary conditions and cell lists in
\texttt{domain.py} (all under \texttt{src/rectsim/}).

\subsection{Vicsek-style noisy alignment}

We consider a discrete-time Vicsek-type alignment model for self-driven agents. At discrete times
$t_n=n\Delta t$ (\cv{dt}), agent $i$ has position $x_i^n\in\Omega\subset\mathbb{R}^2$ (\cv{traj}) and unit heading
$p_i^n := (\cos\theta_i^n,\ \sin\theta_i^n)\in\mathbb{S}^1$, with constant-speed velocity $v_i^n=v_0\,p_i^n$ for a fixed $v_0>0$ (\cv{speed}).
Given an interaction radius $R>0$ (\cv{R}), the neighbor set is $\mathcal{N}_i^n := \{\, j\in\{1,\dots,N\} : d_\Omega(x_i^n,x_j^n) \le R \,\}$, where $d_\Omega$ denotes the domain distance. The local mean heading direction is obtained by averaging neighbors' headings and renormalizing:
\[
s_i^n := \sum_{j\in\mathcal{N}_i^n} p_j^n,
\qquad
\bar p_i^n :=
\begin{cases}
s_i^n/\|s_i^n\|, & \|s_i^n\|>\varepsilon,\\
p_i^n, & \text{otherwise},
\end{cases}
\]
with a small $\varepsilon>0$ to avoid division by zero.
\noindent Equivalently, one may write $\bar\theta_i^n=\mathrm{Arg}(\bar p_i^n)=\mathrm{atan2}((\bar p_i^n)_2,(\bar p_i^n)_1)$ and update angles modulo $2\pi$. The noisy Vicsek update rotates the mean direction by an angular perturbation $\Delta\theta_i^n$:
\begin{equation}
p_i^{n+1} = R(\Delta\theta_i^n)\,\bar p_i^n,
\qquad
R(\phi)=
\begin{pmatrix}
\cos\phi & -\sin\phi\\
\sin\phi & \cos\phi
\end{pmatrix},
\label{eq:vicsek_heading_vec}
\end{equation}
followed by the forward Euler position update
\begin{equation}
x_i^{n+1} = x_i^n + \Delta t\,v_0\,p_i^{n+1}.
\label{eq:vicsek_position}
\end{equation}
This is the standard Vicsek alignment mechanism \cite{vicsek1995vicsek}.
The noise model is selected by \cv{noise.kind}; the original model uses uniform angular noise
$\Delta\theta_i^n \sim \mathrm{Unif}(-\eta/2,\ \eta/2)$ with strength $\eta$ (\cv{eta}); in addition, we also consider Gaussian noise
$\Delta\theta_i^n \sim \mathcal{N}(0,\sigma^2)$.
When \cv{match\_variance} is set, their variances are matched via $\sigma = \eta/\sqrt{12}$.

%
%
%

\subsection{Boundary conditions}
\label{sec:boundary_conditions}

We consider a rectangular domain $\Omega = [0,L_x)\times[0,L_y)$ (\cv{Lx}, \cv{Ly}),
and impose either \emph{periodic} or \emph{reflecting} boundary conditions (\cv{bc}). In periodic boundaries, distances are computed using the minimal-image convention,
\[
d_\Omega(x,y)
=
\sqrt{
\min\!\big(|x_1-y_1|,\ L_x-|x_1-y_1|\big)^2
+
\min\!\big(|x_2-y_2|,\ L_y-|x_2-y_2|\big)^2
},
\]
and periodicity is enforced componentwise after each position update: $x_{i,\alpha}^{n+1} \leftarrow x_{i,\alpha}^{n+1} \bmod L_\alpha$ for $\alpha\in\{1,2\}$.

In reflecting boundaries: $d_\Omega(x,y)=\|x-y\|_2$.  Reflection is enforced coordinatewise (analogously for the $y$-direction): if $\tilde x_{i,1}^{n+1}<0$ then $x_{i,1}^{n+1}=-\tilde x_{i,1}^{n+1}$ and $p_{i,1}^{n+1}=-p_{i,1}^{n+1}$; if $\tilde x_{i,1}^{n+1}>L_x$ then $x_{i,1}^{n+1}=2L_x-\tilde x_{i,1}^{n+1}$ and $p_{i,1}^{n+1}=-p_{i,1}^{n+1}$.

%
%
%
\subsection{Attraction--repulsion interactions with Morse potentials (D'Orsogna-type)}
\label{sec:morse_forces}

To model swarming, milling, and clustering beyond pure alignment, we incorporate short-range
repulsion and longer-range attraction via a Morse-type interaction force, following the D'Orsogna
model~\cite{dorsogna2006morse} and its computational analysis in~\cite{bhaskar2019topology};
the pairwise force is evaluated by \texttt{morse\_force()} in \texttt{morse.py}.

\emph{Morse interaction force.}\;Let $r_{ij} := x_i - x_j$, $r := \|r_{ij}\|$, and $\hat r_{ij} := r_{ij}/r$. The radial Morse force magnitude is
\begin{equation}
f(r)
=
\frac{C_r}{\ell_r} e^{-r/\ell_r}
-
\frac{C_a}{\ell_a} e^{-r/\ell_a},
\label{eq:morse_fr}
\end{equation}
where $C_r,C_a>0$ (\cv{Cr}, \cv{Ca}) denote repulsion and attraction strengths and $\ell_r,\ell_a>0$ (\cv{lr}, \cv{la}) their  interaction lengths. The force exerted by agent $j$ on $i$ is
\begin{equation}
F_{ij}
=
f(\|x_i-x_j\|)\,\hat r_{ij},
\qquad
F_i
=
\sum_{j\neq i} F_{ij}.
\label{eq:morse_vector_force}
\end{equation}
This force derives from the Morse potential $U(r)
=
C_r e^{-r/\ell_r}
-
C_a e^{-r/\ell_a},
\qquad
F_{ij} = -\nabla_{x_i} U(\|x_i-x_j\|)$. A cohesive regime is typically obtained when $\ell_r<\ell_a$, so that repulsion dominates at short distances while attraction dominates at intermediate distances. In our model we primarily use a \emph{direction-only-constant-speed} mechanism that allows Morse forces to steer directions. Let $p_i^n$ denote the unit heading and $\bar p_i^n$ the normalized alignment direction. We define
\begin{equation}
\tilde p_i^n
=
\bar p_i^n
+
\mu_t\,\frac{F_i^n}{\|F_i^n\|+\varepsilon},
\qquad
p_i^{n+1}
=
\frac{\tilde p_i^n}{\|\tilde p_i^n\|},
\qquad
v_i^{n+1}
=
v_0\,p_i^{n+1},
\label{eq:direction_only_force_update}
\end{equation}
where $\mu_t$ (\cv{mu\_t}) is the translational mobility controlling how strongly forces affect particle motion and $\varepsilon>0$ prevents division by zero, ensuring $\|v_i^{n+1}\| = v_0$.
This corresponds to \cv{speed\_mode}~=~\texttt{constant\_with\_forces} in the YAML configuration.
This regime yields clean latent dynamics while still producing collective behaviors such as milling, clustering, and ring formation. We also allow \emph{variable-speed} force-driven dynamics (\cv{speed\_mode}~=~\texttt{variable}): \begin{equation}
v_i(t+\Delta t)
=
v_i(t)
+
\Delta t\,\mu_t\,F_i(t),
\qquad
x_i(t+\Delta t)
=
x_i(t)
+
\Delta t\,v_i(t+\Delta t),
\end{equation}
\noindent For stability  we enforce $v_0\,\Delta t \;\lesssim\; 0.5\,R$ and $r_{\mathrm{cut}} \;\gtrsim\; 3\max(\ell_r,\ell_a)$, where $r_{\mathrm{cut}} = \texttt{rcut\_factor}\cdot\max(\ell_r,\ell_a)$ (\cv{rcut\_factor}). Forces are computed using symmetric pairwise interactions and linked-cell neighbor lists (\texttt{build\_cells()} in \texttt{domain.py}), with minimal-image displacements in periodic boundary conditions.

\begin{breakablealgorithm}
\caption{Discrete-time Vicsek simulation with steering Morse forces}
\label{alg:vicsek}

{\small
\texttt{%
\# Entry point: simulate\_backend() in vicsek\_discrete.py}\\[2pt]

\begin{algorithmic}[1]
\Require $N,\,L_x,L_y,\,\Delta t,\,T,\,v_0,\,R$, \texttt{bc} $\in\{\texttt{periodic},\texttt{reflecting}\}$,
\texttt{noise.kind} $\in\{\texttt{uniform},\texttt{gaussian}\}$, $\eta$, \texttt{match\_variance}, \texttt{save\_every}, \texttt{neighbor\_rebuild},
\texttt{speed\_mode} $\in\{\texttt{constant},\texttt{constant\_with\_forces},\texttt{variable}\}$,
(optional) \texttt{forces.enabled} and Morse params $(\texttt{Cr},\texttt{Ca},\texttt{lr},\texttt{la},\texttt{rcut\_factor},\texttt{mu\_t})$.
\Ensure Saved frames $\{x^n,p^n\}$ at requested indices; metadata + diagnostics.

\State Initialize positions $x^0\subset [0,L_x)\times[0,L_y)$ and angles $\theta^0\sim \mathrm{Unif}(0,2\pi)$.
\State Set headings $p^0 \gets \texttt{headings\_from\_angles}(\theta^0)$; set $n_{\max} \gets \lfloor T/\Delta t\rfloor$.
\State Set $r_{\mathrm{cut}} \gets \texttt{rcut\_factor}\cdot \max(\texttt{lr},\texttt{la})$ (if forces enabled).
\State Build initial cell list: $\texttt{cell\_list} \gets \texttt{build\_cells}(x^0,L_x,L_y,\max(R,r_{\mathrm{cut}}),\texttt{bc})$.

\For{$n=0,1,\dots,n_{\max}-1$}

  \If{$n \bmod \texttt{neighbor\_rebuild}=0$}
    \State $\texttt{cell\_list} \gets \texttt{build\_cells}(x^n,L_x,L_y,\max(R,r_{\mathrm{cut}}),\texttt{bc})$.
  \EndIf

  \State $\mathcal{N}^n \gets \texttt{compute\_neighbors}(x^n,L_x,L_y,R,\texttt{bc},\texttt{cell\_list})$ \Comment{includes self}

  \If{\texttt{forces.enabled}}
    \State $F^n \gets \texttt{morse\_force}(x^n,L_x,L_y,\texttt{bc},\texttt{Cr},\texttt{Ca},\texttt{lr},\texttt{la},r_{\mathrm{cut}},\texttt{cell\_list})$.
  \EndIf

  \For{$i=1,\dots,N$}

    \State $s \gets \sum_{j\in\mathcal{N}_i^n} p_j^n$.
    \State $\bar p_i^n \gets s/\|s\|$ \Comment{alignment mean direction (renormalized)}

    \If{\texttt{speed\_mode} is \texttt{constant\_with\_forces}}
      \State $\widehat F_i^n \gets F_i^n/(\|F_i^n\|+\varepsilon)$.
      \State $\tilde p_i^n \gets \bar p_i^n + \texttt{mu\_t}\,\widehat F_i^n$.
      \State $\bar p_i^n \gets \tilde p_i^n/\|\tilde p_i^n\|$.
    \EndIf

    \State $\phi_i^n \gets \texttt{angle\_noise}(\texttt{rng},\texttt{noise.kind},\eta,\texttt{match\_variance})$.
    \State $p_i^{n+1} \gets \texttt{rotation}(\phi_i^n)\,\bar p_i^n$.
    \State $p_i^{n+1} \gets p_i^{n+1}/\|p_i^{n+1}\|$.

    \If{\texttt{speed\_mode} is \texttt{variable}}
      \State $v_i^{n+1} \gets v_i^n + \Delta t\,\texttt{mu\_t}\,F_i^n$.
      \State $\tilde x_i^{n+1} \gets x_i^n + \Delta t\,v_i^{n+1}$.
    \Else
      \State $\tilde x_i^{n+1} \gets x_i^n + \Delta t\,v_0\,p_i^{n+1}$.
    \EndIf

    \State $(x_i^{n+1},\texttt{flips}_i)\gets \texttt{apply\_bc}(\tilde x_i^{n+1},L_x,L_y,\texttt{bc})$.
    \If{\texttt{bc} is \texttt{reflecting}}
      \State Flip reflected components of $p_i^{n+1}$ according to \texttt{flips}$_i$.
    \EndIf

  \EndFor

  \If{$(n+1)$ is a multiple of \texttt{save\_every}}
    \State Save frame $(x^{n+1},p^{n+1})$.
  \EndIf

\EndFor

\State Store run metadata: $(N,L_x,L_y,\Delta t,T,v_0,R,\texttt{bc},\texttt{noise.kind},\eta,\texttt{speed\_mode})$ and (if enabled)
$(\texttt{Cr},\texttt{Ca},\texttt{lr},\texttt{la},\texttt{mu\_t},\texttt{rcut\_factor})$.
\State Optionally store diagnostics (e.g.\ order parameters computed downstream from saved frames).
\end{algorithmic}
}
\end{breakablealgorithm}

\begin{figure}[H]
\centering

\begin{tabular}{@{}c@{\;}c@{\;}c@{}}
  \feedbackgraphic{0.32\textwidth}{feedback_images/traj_DYN1_t0.png} &
  \feedbackgraphic{0.32\textwidth}{feedback_images/traj_DYN1_t15.png} &
  \feedbackgraphic{0.32\textwidth}{feedback_images/traj_DYN1_t38.png} \\[-2pt]
  \multicolumn{3}{c}{\small (a) Gentle crawl regime (weak Morse forces, low speed)} \\[6pt]

  \feedbackgraphic{0.32\textwidth}{feedback_images/traj_DYN4_t0.png} &
  \feedbackgraphic{0.32\textwidth}{feedback_images/traj_DYN4_t15.png} &
  \feedbackgraphic{0.32\textwidth}{feedback_images/traj_DYN4_t38.png} \\[-2pt]
  \multicolumn{3}{c}{\small (b) Blackhole regime (strong attraction, near-singular collapse)} \\[6pt]

  \feedbackgraphic{0.32\textwidth}{feedback_images/traj_DYN7_t0.png} &
  \feedbackgraphic{0.32\textwidth}{feedback_images/traj_DYN7_t15.png} &
  \feedbackgraphic{0.32\textwidth}{feedback_images/traj_DYN7_t38.png} \\[-2pt]
  \multicolumn{3}{c}{\small (c) Pure Vicsek regime (alignment only, no Morse forces)} \\[2pt]

  {\small $t \approx 0$\,s} & {\small $t \approx 6$\,s} & {\small $t \approx 15$\,s}
\end{tabular}

\caption{Particle trajectories generated by Algorithm~\ref{alg:vicsek} under
  three contrasting parameter regimes, each starting from a uniform initial
  condition.  Arrows indicate heading directions.
  (a)~Gentle crawl: slow drift with gradual clustering from weak Morse
  attraction.
  (b)~Blackhole: rapid collapse into a dense aggregate driven by strong
  attraction ($C_a{=}20$).
  (c)~Pure Vicsek: alignment-only dynamics (no forces) producing coherent
  flocking without spatial aggregation.}
\label{fig:trajectory_snapshots}
\end{figure}

\noindent \emph{System Nonlinearity:} For fixed timestep $t=n$ and given configuration $\mathbf{s}^n = (x^n,p^n)$, the one-step update $\mathbf{s}^{n+1} = T(\mathbf{s}^n)$ is built from smooth operations. However, all of these are applied after computing the neighbor sets $\mathcal{N}_i^n = \{ j : d_\Omega(x_i^n,x_j^n) \le R \},$ The map $x \mapsto \mathcal{N}_i(x)$ is not smooth: when $d_\Omega(x_i,x_j) \le R$, agent $j$ contributes to $i$'s alignment and forces, but when $d_\Omega(x_i,x_j) > R$, it does not. At the threshold $d_\Omega(x_i,x_j) = R$, the contribution switches on or off discontinuously.
The Vicsek update therefore defines a nonlinear one-step map on the microscopic state due to: (i) averaging followed by normalization in the alignment rule, (ii) composition of normalization $v \mapsto v/\|v\|$ with rotational updates $R(\phi)$, (iii) state-dependent discontinuous changes in neighbor sets as particles move, and (iv) nonlinear attraction--repulsion forces $f(r)$ combined with directional normalization. Together, these features produce a high-dimensional, piecewise-smooth dynamical system whose evolution cannot be captured by linear updates.

\subsection{Order parameters for emergent behavior}
\label{sec:order_parameters}

To characterize emergent collective behavior and to validate reduced-order models,
we compute a set of order parameters over time; all are implemented in
\texttt{standard\_metrics.py}.
Closely related diagnostics for attraction--repulsion systems appear in
\cite{bhaskar2019topology}. Let the center of mass and relative positions be defined by
$x_{\mathrm{cm}}^n := \frac{1}{N}\sum_{i=1}^N x_i^n,
\qquad
r_i^n := x_i^n - x_{\mathrm{cm}}^n$,
Global alignment is quantified by the \emph{polarization} order parameter
(\texttt{polarization()}):
\begin{equation}
\Phi^n
:= \left\|\frac{1}{N}\sum_{i=1}^N \frac{v_i^n}{\|v_i^n\|}\right\|
\in [0,1],
\label{eq:polarization}
\end{equation}
\emph{Angular momentum}\,(milling, \texttt{angular\_momentum()}): measures rotational motion about the center of mass
\begin{equation}
M_{\mathrm{ang}}^n
:= \frac{\left|\sum_{i=1}^N r_i^n \times v_i^n\right|}
{N\,\langle\|r\|\rangle\,\langle\|v\|\rangle},
\qquad
\langle\|r\|\rangle := \frac{1}{N}\sum_{i=1}^N \|r_i^n\|,
\quad
\langle\|v\|\rangle := \frac{1}{N}\sum_{i=1}^N \|v_i^n\|.
\label{eq:angular_momentum}
\end{equation}
The \emph{mean speed} (\texttt{mean\_speed()})
\begin{equation}
\langle \|v\| \rangle^n := \frac{1}{N}\sum_{i=1}^N \|v_i^n\|
\label{eq:mean_speed}
\end{equation}
serves as a diagnostic for numerical stability.  \emph{Nematic order} (axial alignment, \texttt{nematic\_order()})
Axial, head--tail symmetric alignment is quantified by the nematic tensor
\begin{equation}
Q^n := \frac{1}{N}\sum_{i=1}^N \left(n_i^n \otimes n_i^n\right) - \frac{1}{2}I,
\qquad
n_i^n := \frac{v_i^n}{\|v_i^n\|},
\qquad
S^n := \lambda_{\max}(Q^n)
\label{eq:nematic_tensor}
\end{equation}
where $\lambda_{\max}$ denotes the largest eigenvalue. Let $\rho^n \in \mathbb{R}^{N_x\times N_y}$ denote the discretized density field at time $t_n$.
We define the \emph{spatial order} parameter (density heterogeneity, \texttt{spatial\_order()}) as
\begin{equation}
\mathcal{O}_{\mathrm{spatial}}^n
:= \operatorname{std}\!\left(\rho^n\right)
= \sqrt{\frac{1}{N_x N_y}
\sum_{a,b}\left(\rho_{a,b}^n - \bar\rho^n\right)^2},
\qquad
\bar\rho^n := \frac{1}{N_x N_y}\sum_{a,b}\rho_{a,b}^n.
\label{eq:spatial_order}
\end{equation}
Large values indicate clustered or heterogeneous spatial structure,
while values near zero correspond to near-uniform density distributions.
This diagnostic is used to compare true and predicted density fields in ROM evaluations.


%
%
%
%
%

\section{Kernel density estimation and bandwidth selection}
\label{subsec:kde}
Following the equation-free framework of \cite{alvarez2025eqfree}, this step serves as a coarse-graining  operator that maps discrete positions $\{x_i(t)\}_{i=1}^N$ (\cv{traj}) to a smooth \emph{number density} field $\rho(\cdot,t)$ (\cv{rho}).
The density movie is produced by \texttt{kde\_density\_movie()} in \texttt{legacy\_functions.py}.
In continuous form, we define
\begin{equation}
\rho(x,t)
=
\sum_{i=1}^N K_h\!\bigl(x-x_i(t)\bigr),
\label{eq:kde_def}
\end{equation}
where $K_h$ is a smooth, nonnegative kernel with bandwidth $h$ (\cv{bandwidth}, specified in grid-cell units so $h=\sigma\,\Delta x$) controlling the smoothing level and normalized so that
$\int_{\mathbb{R}^2} K_h(z)\,dz = 1$.
With this convention, $\rho$ has units of \emph{particles per unit area} and satisfies the mass constraint
$\int_\Omega \rho(x,t)\,dx = N$ (up to discretization), which is essential for downstream POD-based compression and mass-preserving lifting.

We use an isotropic Gaussian kernel,
\begin{equation}
K_h(z)
=
\frac{1}{2\pi h^2}
\exp\!\left(-\frac{\|z\|^2}{2h^2}\right),
\label{eq:gaussian_kernel}
\end{equation}
which provides smooth coarse-graining appropriate for periodic rectangular domains.

\begin{claim}[Histogram plus Gaussian smoothing approximates gridded Gaussian KDE]
\label{claim:histogram_kde}
Let $\tilde\rho(\cdot,t)$ be the piecewise-constant histogram number density on a uniform grid with spacings
$\Delta x = L_x/n_x$ and $\Delta y = L_y/n_y$.
Let $G_\sigma$ denote the discrete Gaussian convolution operator (standard deviation $\sigma$ in \emph{grid-cell units}; \texttt{scipy.ndimage.gaussian\_filter}) applied to the gridded field $\tilde\rho$.
If the grid is square ($\Delta x=\Delta y$) and the same $\sigma$ is used in both axes, then the smoothed histogram
\begin{equation}
\rho(\cdot,t) \;=\; G_\sigma * \tilde\rho(\cdot,t)
\label{eq:smoothed_histogram}
\end{equation}
approximates the \emph{cell-averaged} continuous Gaussian KDE \eqref{eq:kde_def} with physical bandwidth $h=\sigma\,\Delta x$.
\end{claim}

\begin{proof}
Define the histogram number density
\begin{equation}
\tilde\rho(x,t)
=
\frac{1}{\Delta x\,\Delta y}\sum_{i=1}^N \mathbf{1}_{\mathrm{cell}(x_i(t))}(x),
\label{eq:histogram_field}
\end{equation}
so that $\int_\Omega \tilde\rho(x,t)\,dx = N$ holds exactly.
Applying a normalized Gaussian smoothing operator (implemented as a discrete convolution on the grid) yields a gridded field that is well-approximated by convolving $\tilde\rho$ against a Gaussian of physical standard deviation $h=\sigma\,\Delta x$:
\begin{equation}
(G_\sigma * \tilde\rho)(x,t)
\approx
\int_\Omega \tilde\rho(x',t)\,
\frac{1}{2\pi h^2}\exp\!\left(-\frac{\|x-x'\|^2}{2h^2}\right)\,dx'.
\label{eq:gaussian_convolution}
\end{equation}
Substituting \eqref{eq:histogram_field} into \eqref{eq:gaussian_convolution} gives
\begin{equation}
(G_\sigma * \tilde\rho)(x,t)
\approx
\sum_{i=1}^N \frac{1}{\Delta x\,\Delta y}
\int_{\mathrm{cell}(x_i(t))}
\frac{1}{2\pi h^2}\exp\!\left(-\frac{\|x-x'\|^2}{2h^2}\right)\,dx'.
\label{eq:cell_avg_kde}
\end{equation}
The right-hand side is precisely a \emph{cell-averaged} Gaussian KDE: each particle contributes a Gaussian kernel integrated over its occupied cell and normalized by the cell area.
As $\Delta x,\Delta y\to 0$ (for fixed $h$), the cell averages converge to pointwise kernel evaluations and \eqref{eq:cell_avg_kde} approaches $\sum_{i=1}^N K_h(x-x_i(t))$.
Thus, histogram deposition followed by Gaussian smoothing produces a consistent gridded approximation to the Gaussian KDE \eqref{eq:kde_def}.
\end{proof}

\begin{remark}[Mass conservation in periodic domains]
\label{rem:kde_mass}
For periodic boundary conditions, we apply Gaussian smoothing with \texttt{mode="wrap"}.
Since the discrete Gaussian filter is normalized (weights sum to~$1$), discrete convolution preserves the total discrete mass: $\sum_{a,b} (G_\sigma * \tilde\rho)_{ab} = \sum_{a,b} \tilde\rho_{ab}$. Therefore the total mass $M(t) := \sum_{a,b} \rho_{ab}(t)\,\Delta x\,\Delta y$ is preserved up to floating-point error and equals $N$ for number density (or~$1$ for probability density normalization).
\end{remark}

\begin{breakablealgorithm}
\caption{Fast gridded Gaussian KDE (histogram + Gaussian smoothing)}
\label{alg:kde_hist_gauss}
\texttt{%
\# Entry point: kde\_density\_movie() in legacy\_functions.py}\\[2pt]
\begin{algorithmic}[1]
\Require Trajectory positions $\mathbf{x}(t)\in\mathbb{R}^{N\times 2}$ at time $t$,
domain size $(L_x,L_y)$, grid resolution $(n_x,n_y)$,
bandwidth $\sigma$ (in grid cells), boundary condition \texttt{bc} $\in\{\texttt{"periodic"},\texttt{"nonperiodic"}\}$
\Ensure Density field $\rho(t)\in\mathbb{R}^{n_y\times n_x}$ (row-major: $y$ index first)
\State $\Delta x\gets L_x/n_x$, $\Delta y\gets L_y/n_y$
\State Extract arrays: $x\gets\mathbf{x}[:,0]$, $y\gets\mathbf{x}[:,1]$
\If{\texttt{bc} $=$ \texttt{"periodic"}}
  \State Wrap positions: $x\gets x\bmod L_x$, $y\gets y\bmod L_y$
\EndIf
\State Compute 2D histogram $H(t)\in\mathbb{R}^{n_y\times n_x}$ from $(y,x)$ using bin edges on $[0,L_y]\times[0,L_x]$
\Comment{Axis convention check: a single particle at $(x,y)$ increments $H[\lfloor y/\Delta y\rfloor,\lfloor x/\Delta x\rfloor]$}
\State Convert counts to number density: $\tilde\rho(t)\gets H(t)/(\Delta x\,\Delta y)$
\If{\texttt{bc} $=$ \texttt{"periodic"}}
  \State Smooth with periodic wrap: $\rho(t)\gets G_\sigma(\tilde\rho(t))$ using \texttt{mode="wrap"}
\Else
  \State Smooth with edge extension: $\rho(t)\gets G_\sigma(\tilde\rho(t))$ using \texttt{mode="nearest"}
  \Comment{Note: \texttt{nearest} repeats edge values; this is not reflective padding.}
\EndIf
\State Compute mass: $M(t)\gets\sum_{a,b}\rho_{ab}(t)\,\Delta x\,\Delta y$
\State Compute relative drift: $\varepsilon(t)\gets |M(t)-N|/N$
\If{$\varepsilon(t) > 10^{-3}$}
  \State \textbf{print} warning: mass deviation $\varepsilon(t)$ detected (typically indicates wrapping/binning errors)
\EndIf
\State \Return $\rho(t)$ and metadata $\{\sigma,n_x,n_y,L_x,L_y,\texttt{bc}\}$
\end{algorithmic}
\end{breakablealgorithm}

\emph{Implementation details:} We specify the bandwidth $\sigma$ in grid-cell units, so the physical smoothing length is $h=\sigma\,\Delta x$ (for a square grid $\Delta x=\Delta y$).
Typical production values are $\sigma\in[2.0,3.0]$ with $n_x=n_y=64$.
We record metadata (bandwidth, grid resolution, domain size, boundary condition) alongside the resulting density movie $\rho(t)\in\mathbb{R}^{T\times n_y\times n_x}$, which serves as input to POD-based compression and subsequent reduced-order modeling.

\subsection{Multi-field coarse-graining}
\label{sec:multifield_cg}

Beyond the scalar density $\rho$, the multi-field WSINDy formulation (Chapter~6) requires
two additional coarse-grained fields and, optionally, a nonlocal interaction field.
All three are computed by routines in \texttt{fields.py} using the same
histogram-plus-Gaussian-smoothing approach as the density KDE above.

\paragraph{Polarisation flux fields.}
The velocity-weighted \emph{polarisation flux} fields $p_x,p_y$ (\cv{px}, \cv{py})
are defined by replacing the unit weight in~\eqref{eq:kde_def} with the
corresponding velocity component:
\begin{equation}
p_\alpha(x,t)
=
\sum_{i=1}^N K_h\!\bigl(x-x_i(t)\bigr)\,v_{i,\alpha}(t),
\qquad \alpha\in\{x,y\},
\label{eq:flux_kde}
\end{equation}
where $K_h$ is the same Gaussian kernel used for $\rho$.
In practice, \texttt{compute\_flux\_kde()} deposits velocity-weighted counts
into a 2-D histogram and applies \texttt{gaussian\_filter} with the same
bandwidth $\sigma$ and periodic-wrap mode, yielding
$p_x,p_y\in\mathbb{R}^{T\times n_y\times n_x}$.

\paragraph{Morse potential field.}
The Morse interaction potential field $\Phi$ (\cv{Phi}) is the convolution
of the density with the radial Morse kernel
$W(r) = -C_r\,e^{-r/\ell_r} + C_a\,e^{-r/\ell_a}$ (\cv{W}):
\begin{equation}
\Phi(x,t)
=
(W * \rho)(x,t)
=
\mathcal{F}^{-1}\!\bigl[\hat W \cdot \hat\rho\bigr](x,t),
\label{eq:morse_potential_field}
\end{equation}
evaluated efficiently via 2-D FFT convolution by
\texttt{compute\_morse\_potential()} in \texttt{fields.py}.
The combined field set $(\rho, p_x, p_y)$---together with their spatial
derivatives and, when forces are active, $\Phi$---forms the input to the
weak-form PDE library constructed in Chapter~6.
 
%
\chapter{Experimental Protocol and Metrics}
\label{ch:exp_protocol_metrics}

This chapter describes the evaluation protocol shared by every experiment in
this thesis.  Section~\ref{sec:pipeline_orchestration} summarises the eight-stage
pipeline that takes a YAML configuration file through to forecast metrics.
Section~\ref{sec:splits} defines the train/test split.
Section~\ref{sec:lag_selection} describes the information-theoretic lag
selection study that motivates the MVAR lag order.
Section~\ref{sec:metrics} defines all evaluation metrics and the forecast
horizon.

%
\section{Pipeline orchestration}
\label{sec:pipeline_orchestration}

Every experiment is driven by a single YAML configuration file (one per
regime) and executed end-to-end by
\cv{ROM\_WSINDY\_pipeline.py}.
The pipeline proceeds through the eight sequential stages shown in
Figure~\ref{fig:pipeline_stages}.

\begin{figure}[h]
\centering
\begin{tikzpicture}[
    node distance=0.45cm and 0.0cm,
    stage/.style={
      draw, rounded corners=3pt, minimum width=12.5cm,
      minimum height=0.72cm, font=\small, align=center,
      fill=blue!6
    },
    arr/.style={-{Stealth[length=4pt]}, thick}
  ]
  \node[stage]                  (s1) {\textbf{1.}~Generate training data --- 168 Vicsek--Morse trajectories, 4~IC families};
  \node[stage, below=of s1]     (s2) {\textbf{2.}~Build shared POD basis --- shift-align, $\sqrt{\cdot}$-transform, truncated SVD $\to$ $r$ modes};
  \node[stage, below=of s2]     (s3) {\textbf{3.}~Build latent dataset --- project $\to$ $y(t)\in\mathbb{R}^r$, subsample$\,{=}\,3$, construct $(X,Y)$ pairs};
  \node[stage, below=of s3]     (s4) {\textbf{4.}~Train MVAR \& LSTM on shared latent dataset (Chapter~\ref{ch:latent_rom})};
  \node[stage, below=of s4]     (s5) {\textbf{5.}~Generate test data --- 4 held-out trajectories (one per IC family)};
  \node[stage, below=of s5]     (s6) {\textbf{6.}~Evaluate MVAR \& LSTM --- condition, forecast, lift, compute metrics};
  \node[stage, below=of s6]  (s7) {\textbf{7.}~WSINDy PDE discovery --- same train split, raw-field weak systems, MSTLS};
  \node[stage, below=of s7]  (s8) {\textbf{8.}~WSINDy test evaluation --- ETDRK4 forecast \& same metric suite};
  %
  \draw[arr] (s1) -- (s2);
  \draw[arr] (s2) -- (s3);
  \draw[arr] (s3) -- (s4);
  \draw[arr] (s4) -- (s5);
  \draw[arr] (s5) -- (s6);
  \draw[arr] (s6) -- (s7);
  \draw[arr] (s7) -- (s8);
\end{tikzpicture}
\caption{Eight-stage pipeline orchestrated by \cv{ROM\_WSINDY\_pipeline.py}.
  Stages~1--6 run the EF-ROM branch;
  stages~7--8 run WSINDy PDE discovery and evaluation.}
\label{fig:pipeline_stages}
\end{figure}

\noindent
All results, per-timestep metric tables, and model artefacts are written to a
self-contained experiment directory.  A summary
file \cv{summary.json} records the mean test $R^2$, training times, and
(for WSINDy) the discovered equation and weak-form $R^2$.  The script
\cv{rerun\_evaluation.py} can re-evaluate stored models on new test data or
with a different forecast start without retraining.

\noindent
Stages~1 and~5 define a common upstream dataset for all models: the same
particle simulations, the same KDE-based coarse-grained fields, and the same
train/test trajectory split.  The branches diverge only after coarse-graining.
Stages~2--6 build a POD-based latent dataset and learn discrete-time latent
evolution maps for MVAR and LSTM, whereas stages~7--8 bypass POD and assemble
weak-form regression systems directly from raw density and auxiliary fields on
the same training trajectories.

%
\section{Train/test splits}
\label{sec:splits}

Training and test data are drawn from the same parameter regime but use
disjoint initial conditions.

\paragraph{Training set.}
Each regime generates 168 runs from four IC families
(Table~\ref{tab:ic_split}, Chapter~\ref{ch:experimentation}).
For the EF-ROM branch, all runs contribute to the shared POD basis and latent
dataset.  For WSINDy, the same training trajectories are used in physical
space: density and auxiliary fields are computed on each run and then
converted into weak-form regression constraints.

\paragraph{Test set.}
Four runs---one per IC family---are generated at held-out IC parameters
(e.g.\ a Gaussian centred at $(12.5,12.5)$ rather than the training grid
positions) and simulated for $T_{\mathrm{test}}$ (Table~\ref{tab:fixed_params}).
They are never seen during POD construction, latent-model fitting, or WSINDy
regression.  The same held-out trajectories define the test set for both
branches.

\paragraph{Forecasting protocol.}
Evaluation is performed in closed-loop (iterated single-step) forecasting.
Let $t_{\mathrm{fc}}^{\mathrm{req}}$ denote the requested
\cv{forecast\_start}, and let
$w_{\max}=\max\{w_m:\text{$m$ enabled ROM model in the config}\}$ be the
largest enabled latent lag.  The effective common forecast start is then
\begin{equation}
  t_{\mathrm{fc}}
  \;=\;
  \max\!\bigl(t_{\mathrm{fc}}^{\mathrm{req}},\, w_{\max}\,\Delta t_{\mathrm{rom}}\bigr).
\end{equation}
All ROM models therefore begin autoregressive rollout from the same
physical time, while each still receives a full ground-truth conditioning
window.  In the baseline systematic configs,
$w_{\mathrm{MVAR}}=w_{\mathrm{LSTM}}=5$, so
$t_{\mathrm{fc}}=0.6$\,s at $\Delta t_{\mathrm{rom}}=0.12$\,s; in
lag-ablation configs with larger LSTM windows, the common start shifts
later accordingly.  The same effective forecast-start time is also used
for WSINDy evaluation, but there the model is initialised from physical
fields at $t_{\mathrm{fc}}$ rather than from a lagged latent window.

%
\section{Lag selection}
\label{sec:lag_selection}

The MVAR lag order $w$ determines how many past latent frames the model
uses for prediction.  To choose $w$ in a principled way, we run an
information-theoretic lag selection study on every regime using the script
\cv{alvarez\_lag\_selection.py}.

Given latent trajectories from the training set, we fit unrestricted
VAR($w$) models for $w=1,2,\dots,w_{\max}$ (with $w_{\max}=100$) using
\texttt{statsmodels}' \cv{VAR.select\_order()} interface.  For each lag
order, three information criteria are evaluated (with $k_w = w\,r^2$ parameters and $T$ effective samples): the \textbf{BIC} (Bayesian), $\mathrm{BIC}(w) = \ln|\hat\Sigma_w| + k_w \ln T / T$; the \textbf{HQIC} (Hannan--Quinn), replacing $\ln T$ with $2\ln\ln T$; and the \textbf{AIC} (Akaike), $\mathrm{AIC}(w) = \ln|\hat\Sigma_w| + 2k_w/T$.
The lag that minimises each criterion is recorded.  For the representative
CS01 regime, for example, BIC selects $w^*_{\mathrm{BIC}}=23$,
HQIC selects $w^*_{\mathrm{HQIC}}=37$, and AIC selects
$w^*_{\mathrm{AIC}}=50$.

\paragraph{Empirical lag sweep.}
Because the information criteria assume an unrestricted VAR (every latent
dimension regressed on every other), and our production MVAR uses
independent ridge-regularised fits, we additionally run an empirical lag
sweep to verify that forecast quality is not sensitive to the exact choice
of~$w$.  For the CS01 regime, seven configs
(\cv{configs/lag\_sweep/LSWEEP\_CS01\_*.yaml}) fix the MVAR lag at
$w_{\mathrm{MVAR}}=5$ and vary the LSTM lag over
$w_{\mathrm{LSTM}}\in\{5,10,15,20,25,30,40\}$.
Each config runs the full pipeline end-to-end and records test~$R^2$.

\paragraph{Outcome.}
The sweep confirms that MVAR at $w{=}5$ already achieves strong forecast
performance ($R^2>0.9$ on most regimes), well below the
information-theoretically optimal lag.  This is consistent with the low
effective dimensionality of the aligned latent dynamics: after shift-alignment
and the $\sqrt{\cdot}$-transform, the latent trajectories are
nearly linear, so a compact lag window suffices.  We therefore fix
$w_{\mathrm{MVAR}}=5$ in the baseline systematic comparison.  The LSTM
lag is treated as a separate hyperparameter in dedicated lag ablations;
whenever $w_{\mathrm{LSTM}}\neq w_{\mathrm{MVAR}}$, the common forecast
start is shifted according to the largest enabled lag.

%
\section{Evaluation metrics}
\label{sec:metrics}

All metrics compare predicted densities $\hat{\boldsymbol{\rho}}^n$ against
ground truth $\boldsymbol{\rho}^n$ over the forecast window
$\mathcal{T}_{\mathrm{fc}} = \{n_0, n_0{+}1, \dots, n_1\}$, where $n_0$
is the forecast start index and $n_1$ is the terminal index.
Let $d = N_x N_y$ denote the number of spatial grid points.

\paragraph{Coefficient of determination ($R^2$).}
\label{subsec:r2_time}
The per-snapshot $R^2$ is computed as implemented in
\cv{compute\_r2\_timeseries()}:
\begin{equation}
  R^2(n) \;=\; 1 \;-\;
  \frac{\sum_{j=1}^{d}\bigl(\rho_j^n - \hat\rho_j^n\bigr)^2}
       {\sum_{j=1}^{d}\bigl(\rho_j^n - \bar\rho^n\bigr)^2},
  \qquad
  \bar\rho^n = \frac{1}{d}\sum_{j=1}^d \rho_j^n,
  \label{eq:r2_snapshot}
\end{equation}
where the denominator uses the spatial mean of the ground-truth snapshot at
time~$n$.  We report the mean over the forecast window,
$R^2_{\mathrm{fc}} = |\mathcal{T}_{\mathrm{fc}}|^{-1}\sum_{n\in\mathcal{T}_{\mathrm{fc}}} R^2(n)$,
as the primary scalar metric.

\paragraph{Relative $\ell^p$ errors.}
\label{subsec:rel_errors}
Three relative error norms are computed per snapshot by
\cv{compute\_error\_timeseries()} and averaged over the
forecast window:
\begin{align}
  \mathrm{RelErr}_1(n) &= \frac{\|\boldsymbol{e}^n\|_1}{\|\boldsymbol{\rho}^n\|_1 + \epsilon},
  &
  \mathrm{RelErr}_2(n) &= \frac{\|\boldsymbol{e}^n\|_2}{\|\boldsymbol{\rho}^n\|_2 + \epsilon},
  &
  \mathrm{RelErr}_\infty(n) &= \frac{\|\boldsymbol{e}^n\|_\infty}{\|\boldsymbol{\rho}^n\|_\infty + \epsilon},
  \label{eq:rel_errors}
\end{align}
where $\boldsymbol{e}^n = \boldsymbol{\rho}^n - \hat{\boldsymbol{\rho}}^n$
and $\epsilon = 10^{-12}$ guards against division by zero.

\paragraph{Mass conservation.}
\label{subsec:mass_metric}
The discrete total mass at time $n$ is
$m^n = (\Delta x\,\Delta y)\,\mathbf{1}^\top\boldsymbol{\rho}^n$ (and
analogously $\hat m^n$ for the prediction).  We track $m^n$ and $\hat m^n$
over the forecast window and report the maximum relative mass violation:
\begin{equation}
  \Delta m_{\max} \;=\;\max_{n\in\mathcal{T}_{\mathrm{fc}}}
    \frac{|\hat m^n - m^n|}{m^{n_0}}.
  \label{eq:mass_violation}
\end{equation}
The pipeline optionally applies a post-hoc mass post-processing step
(\cv{mass\_postprocess}), either the global rescaling
\eqref{eq:mass_postprocess_scale} or the simplex projection
\eqref{eq:mass_postprocess_simplex}, to match the initial mass.  In the
preprocessing ablations this had little effect on MVAR but improved LSTM
rollout robustness and reduced $\Delta m_{\max}$ in practice.

\paragraph{Spatial order parameter.}
\label{subsec:spatial_order_metric}
To assess whether the model preserves the spatial structure of the density
field, we track the spatial standard deviation
$\sigma_\rho(n) = \mathrm{std}(\boldsymbol{\rho}^n)$ defined in
Eq.~\eqref{eq:spatial_order} (Chapter~\ref{ch:background}).
A high $\sigma_\rho$ indicates concentrated density; a low value
indicates a nearly uniform field.

\paragraph{Parameter counts.}
\label{subsec:param_counts}
At the baseline comparison settings
($r{=}19$, $w_{\mathrm{MVAR}}{=}w_{\mathrm{LSTM}}{=}5$, hidden width
$h{=}128$, $L{=}2$ LSTM layers), the MVAR has
$p_{\mathrm{MVAR}} = w\,r^2 + r = 1{,}824$ free parameters, while the
LSTM has $p_{\mathrm{LSTM}} \approx 178{,}835$---a roughly $98\times$
ratio.  Full derivations appear in
Chapter~\ref{ch:latent_rom}, \S\ref{subsec:mvar_identifiability}
and~\S\ref{sec:lstm_surrogate}.
 
%
\chapter{Experimentation}
\label{ch:experimentation}

\section{Parameter choices and dynamics regimes}
\label{sec:params_fixed}

To keep comparisons across learning models meaningful, we fix a core set of
numerical and grid parameters across the full pipeline
(Table~\ref{tab:fixed_params}) and vary only the physics and stochastic knobs
that define each dynamics regime.  All configuration files live in
\cv{configs/systematic/}.

\begin{table}[ht!]
\centering
\caption{Fixed parameters shared by all dynamics regimes and pipeline stages.}
\label{tab:fixed_params}
\begin{tabular}{lll}
\hline
Parameter & Value & Code key \\
\hline
Number of particles $N$                    & 300               & \cv{N} \\
Domain size $L_x = L_y$                    & 25                & \cv{Lx}, \cv{Ly} \\
Boundary conditions                        & periodic          & \cv{bc} \\
Time step $\Delta t$                        & 0.04\,s           & \cv{dt} \\
Training horizon $T_{\mathrm{train}}$       & 20\,s (500 steps)  & \cv{T} \\
Test horizon $T_{\mathrm{test}}$            & 200\,s            & \cv{test\_sim.T} \\
Interaction radius $R$                      & 4.0               & \cv{R} \\
KDE grid resolution $n_x \times n_y$       & $64 \times 64$    & \cv{nx}, \cv{ny} \\
KDE bandwidth                               & 5.0               & \cv{bandwidth} \\
Snapshot subsampling stride                 & 3                 & \cv{subsample} \\
Fixed POD modes $m$                         & 19                & \cv{fixed\_modes} \\
Density transform                           & $\sqrt{\cdot}$    & \cv{density\_transform} \\
Shift alignment                             & enabled (mean ref) & \cv{shift\_align} \\
Mass post-processing                        & scale             & \cv{mass\_postprocess} \\
\hline
\end{tabular}
\end{table}

\begin{figure}[H]
\centering
\feedbackgraphic{0.85\textwidth}{feedback_images/config_legend.pdf}
\caption{Colour and marker legend for all dynamics configurations used
  throughout this thesis.  Each symbol identifies a single regime in subsequent
  result figures.}
\label{fig:config_legend}
\end{figure}

\subsection{NDYN regimes (physics-knob variations)}
\label{subsec:ndyn_regimes}

Fourteen \emph{NDYN} regimes isolate individual physics knobs---speed, noise
strength, attraction/repulsion amplitude, force length scale, and speed
mode---while holding the others at moderate baseline values.  All use Morse
forces (\cv{forces.enabled: true}) except NDYN08, which disables forces
entirely (pure Vicsek alignment).  Table~\ref{tab:ndyn_regimes} lists the
varying parameters; all other parameters are as in Table~\ref{tab:fixed_params}.

\begin{table}[ht!]
\centering
\caption{NDYN regimes.  Speed mode is \cv{constant\_with\_forces} unless
  noted.  Morse torque $\mu_t{=}0.3$ and \cv{rcut\_factor}${}=5$ are shared.}
\label{tab:ndyn_regimes}
\scriptsize
\begin{tabular}{llrrrrrrl}
\hline
Label & Nickname & $v_0$ & $\eta$ & $C_a$ & $C_r$ & $\lambda_a$ & $\lambda_r$ & Notes \\
\hline
NDYN01 & crawl       & 0.5 & 0.03 & 0.8  & 0.3  & 1.5 & 0.5 & easy POD baseline \\
NDYN02 & flock       & 3.0 & 0.1  & 2.0  & 0.5  & 1.5 & 0.5 & coherent cluster \\
NDYN03 & sprint      & 10.0& 0.1  & 2.0  & 0.5  & 1.5 & 0.5 & fast transport \\
NDYN04 & gas         & 3.0 & 1.5  & 0.5  & 0.2  & 1.5 & 0.5 & noise destroys structure \\
NDYN05 & blackhole   & 4.0 & 0.15 & 20.0 & 0.05 & 1.5 & 0.3 & near-singular spike \\
NDYN06 & supernova   & 3.0 & 0.15 & 0.05 & 15.0 & 1.5 & 0.5 & explosive repulsion \\
NDYN07 & crystal     & 2.0 & 0.1  & 3.0  & 3.0  & 2.0 & 0.5 & balanced equilibrium \\
NDYN08 & pure Vicsek & 3.0 & 0.3  & ---  & ---  & --- & --- & forces \emph{off} \\
NDYN09 & longrange   & 2.0 & 0.2  & 1.5  & 0.6  & 6.0 & 3.0 & global collective modes \\
NDYN10 & shortrange  & 3.0 & 0.15 & 4.0  & 2.0  & 0.3 & 0.1 & fragmented clusters \\
NDYN11 & noisy collapse & 3.0 & 0.8 & 12.0 & 0.1 & 1.5 & 0.5 & intermittent clusters \\
NDYN12 & fast explosion & 8.0 & 0.2 & 0.1 & 10.0 & 1.5 & 0.5 & rapid dispersal \\
NDYN13 & chaos       & 12.0& 1.2  & 1.0  & 0.5  & 1.5 & 0.5 & ROM lower bound \\
NDYN14 & varspeed    & 3.0 & 0.2  & 2.0  & 0.8  & 1.5 & 0.5 & \cv{variable} speed \\
\hline
\end{tabular}
\end{table}

\subsection{DO regimes (outcome-targeted variations)}
\label{subsec:do_regimes}

A second set of fourteen \emph{DO} (dynamics-outcome) regimes holds speed
$v_0{=}3$, noise $\eta{=}0.2$, and attraction coefficient $C_a{=}1$ fixed, and
varies only the repulsion coefficient $C_r$ and the repulsion length scale
$\lambda_r$ (with $\lambda_a{=}1$ throughout).  This systematic sweep targets
specific collective behaviours observed in D'Orsogna-type
models~\cite{dorsogna2006morse}: collective swarms (CS), single mills (SM), double
mills (DM), double rings (DR), and several escape patterns (EC, ES, EU).
Table~\ref{tab:do_regimes} lists the 14~configurations.

\begin{table}[ht!]
\centering
\caption{DO regimes.  All share $v_0{=}3$, $\eta{=}0.2$, $C_a{=}1$,
  $\lambda_a{=}1$; only $C_r$ and $\lambda_r$ vary.}
\label{tab:do_regimes}
\footnotesize
\begin{tabular}{llrrl}
\hline
Label & Behaviour & $C_r$ & $\lambda_r$ & Notes \\
\hline
CS01 & collective swarm    & 0.1 & 0.5 & \\
CS02 & collective swarm    & 0.5 & 3.0 & \\
CS03 & collective swarm    & 0.9 & 3.0 & \\
SM01 & single mill         & 0.5 & 0.1 & \\
SM02 & single mill         & 3.0 & 0.1 & \\
SM03 & single mill         & 2.0 & 0.5 & \\
DM01 & double mill         & 0.9 & 0.5 & showcase regime \\
DR01 & double ring         & 0.1 & 0.1 & \\
DR02 & double ring         & 0.9 & 0.9 & \\
EC01 & escape (collective) & 2.0 & 3.0 & \\
EC02 & escape (collective) & 3.0 & 0.5 & \\
ES01 & escape (symmetric)  & 3.0 & 0.9 & \\
EU01 & escape (unsymmetric)& 2.0 & 2.0 & \\
EU02 & escape (unsymmetric)& 3.0 & 3.0 & \\
\hline
\end{tabular}
\end{table}

\paragraph{Variable-speed variants.}
\label{subsec:vs_variants}
Most NDYN and DO configurations have a companion \emph{\_VS} variant
(e.g.\ \cv{NDYN01\_crawl\_VS.yaml}) that changes exactly one parameter:
\cv{speed\_mode} from \cv{constant\_with\_forces} to \cv{variable}.  In
variable-speed mode, Morse forces modulate agent \emph{speed} as well as
heading, producing non-constant flux fields.  NDYN08 (pure Vicsek, no forces)
and NDYN14 (already variable-speed) have no \_VS counterpart.

%
\section{Initial condition families}
\label{sec:ic_families}

To probe ROM robustness beyond a single ``typical'' microscopic state, we
generate ensembles of trajectories from four qualitatively different
\emph{initial position} distributions, implemented in
\cv{src/rectsim/initial\_conditions.py}.  The intent is to
(i)~span a range of macroscopic organisations (homogeneous, localised,
symmetric, multi-modal) and (ii)~induce distinct transient regimes under the
\emph{same} simulator and numerical settings.  All particles evolve on the
periodic domain $\Omega=[0,L_x]\times[0,L_y]$.

\paragraph{Common velocity initialisation.}
Across all IC families we draw random headings so that differences in dynamics
are driven by the initial \emph{spatial} configuration:
\begin{equation}
  \theta_i \sim \mathcal{U}(0,2\pi),\qquad
  \mathbf{v}_i(0)=v_0(\cos\theta_i,\sin\theta_i)^\top,\qquad i=1,\dots,N,
  \label{eq:ic_vel_init}
\end{equation}
so $\|\mathbf{v}_i(0)\|=v_0$ for all particles and the initial polarisation is
close to zero in expectation.

\paragraph{Uniform IC.}
\label{subsec:ic_uniform}
The uniform family (\cv{uniform\_distribution}) represents a maximum-entropy
baseline with no imposed spatial structure:
\begin{equation}
  \mathbf{x}_i(0)\sim \mathcal{U}(\Omega), \qquad i=1,\dots,N,
  \label{eq:ic_uniform_def}
\end{equation}
equivalently $x_i\sim\mathcal{U}(0,L_x)$ and $y_i\sim\mathcal{U}(0,L_y)$
independently.  The expected coarse density is spatially constant,
$\mathbb{E}[\rho(\mathbf{r},0)] = N/(L_xL_y)$.  Uniform ICs serve both as
training diversity and as the baseline for generalisation comparisons.

\paragraph{Single Gaussian IC.}
\label{subsec:ic_gaussian}

The single-Gaussian family (\cv{gaussian\_cluster}) produces a localised
``blob'' with controllable centre and spread:
\begin{equation}
  \mathbf{x}_i(0)\sim \mathcal{N}(\boldsymbol{\mu},\sigma^2 I_2),
  \qquad i=1,\dots,N.
  \label{eq:ic_gaussian_def}
\end{equation}
The centre $\boldsymbol{\mu}=(\mu_x,\mu_y)$ is chosen from a $4\times 4$ grid
of interior positions (\cv{positions\_x}, \cv{positions\_y}
$\in\{5,10,15,20\}$) to avoid wrap artefacts at domain boundaries, and
$\sigma^2$ is varied across three values
(\cv{variances} $\in\{1.5,3.0,5.0\}$).  Sampled positions are wrapped back
into~$\Omega$ via modulo.  This yields $4\times 4\times 3=48$ parameter
combinations with \cv{n\_samples\_per\_config}${}=1$, giving
80~training runs (the config rounds to 80 via the enumeration logic).

\paragraph{Ring IC.}
\label{subsec:ic_ring}

The ring family (\cv{ring\_distribution}) enforces rotational symmetry at
$t{=}0$ and typically generates vortex-like macroscopic patterns.  Let
$\mathbf{c}=(L_x/2,\,L_y/2)$ be the domain centre.  Angles and radii are
sampled in polar coordinates:
\begin{equation}
  \theta_i \sim \mathcal{U}(0,2\pi),\qquad
  r_i \sim \mathcal{N}(R_{\mathrm{ring}},\,\sigma_r^2),
  \label{eq:ic_ring_polar}
\end{equation}
and mapped to Cartesian positions
$\mathbf{x}_i(0)=\mathbf{c}+r_i(\cos\theta_i,\sin\theta_i)^\top$, then
wrapped into~$\Omega$ via modulo.  Training enumerates radii
$R_{\mathrm{ring}}\in\{3,5,8\}$ (\cv{radii}) and widths
$\sigma_r\in\{0.5,1.0\}$ (\cv{widths}) with
\cv{n\_samples\_per\_config}${}=4$, yielding $3\times 2\times 4=24$~training
runs per regime.

\paragraph{Two-cluster Gaussian IC.}
\label{subsec:ic_two_clusters}

The two-cluster family (\cv{two\_clusters}) generates a bimodal initial
distribution as an equal-weight mixture of two isotropic Gaussians:
\begin{equation}
  p(\mathbf{x})
  =\tfrac12\,\mathcal{N}(\boldsymbol{\mu}_1,\sigma^2 I_2)
  +\tfrac12\,\mathcal{N}(\boldsymbol{\mu}_2,\sigma^2 I_2),
  \label{eq:ic_two_clusters_mix}
\end{equation}
implemented by sampling $N_1=\lfloor N/2\rfloor$ particles from the first
component and $N_2=N-N_1$ from the second.  The centres are placed
symmetrically about the domain midpoint, separated horizontally by distance~$d$
(\cv{separation}):
\begin{equation}
  \boldsymbol{\mu}_1=\Big(\tfrac{L_x}{2}-\tfrac{d}{2},\,\tfrac{L_y}{2}\Big),
  \qquad
  \boldsymbol{\mu}_2=\Big(\tfrac{L_x}{2}+\tfrac{d}{2},\,\tfrac{L_y}{2}\Big).
  \label{eq:ic_two_clusters_centers}
\end{equation}
Default values are $d=L_x/3$ and $\sigma=\min(L_x,L_y)/8$.  After sampling,
positions are wrapped into~$\Omega$ via modulo.  Training enumerates
$d\in\{4,7,10\}$ and $\sigma\in\{1.5,3.0\}$ with
\cv{n\_samples\_per\_config}${}=4$, yielding $3\times 2\times 4=24$~training
runs per regime.

%
\paragraph{Training--test split summary.}
\label{subsec:train_test_split}

Table~\ref{tab:ic_split} summarises the per-regime dataset composition.  Every
regime (NDYN and DO alike) uses the identical IC schedule, so the total
training set comprises $168 \times 28 = 4{,}704$ simulations across all base
regimes (and the same count again for the \_VS variants that are run).

\begin{table}[H]
\centering
\caption{Training and test runs per dynamics regime.}
\label{tab:ic_split}
\begin{tabular}{lrr}
\hline
IC family & Training runs & Test runs \\
\hline
Gaussian     & 80 & 1 \\
Uniform      & 40 & 1 \\
Ring         & 24 & 1 \\
Two-clusters & 24 & 1 \\
\hline
\textbf{Total} & \textbf{168} & \textbf{4} \\
\hline
\end{tabular}
\end{table}

Each run is assigned a unique \cv{run\_id} from which a deterministic RNG seed
is derived, ensuring exact reproducibility.  The density movies produced by
these simulations feed both modeling branches.  For the EF-ROM pipeline they
are assembled into a global snapshot matrix for POD basis construction, as
described in Chapter~\ref{ch:global_pod_pipeline}.  For WSINDy, the same
regime-specific training and test simulations are retained in physical space
and converted into density and auxiliary multi-fields for weak-form regression.
 
%
\chapter{Global POD Reduction and Restriction--Lifting Pipeline}
\label{ch:global_pod_pipeline}

This chapter describes the data-driven reduction pipeline that maps
high-dimensional density fields to low-dimensional latent trajectories
suitable for time-series modelling.  The pipeline is implemented in
\cv{src/rectsim/pod\_builder.py} (training) and
\cv{src/rectsim/test\_evaluator.py} (evaluation), orchestrated by
\cv{ROM\_WSINDY\_pipeline.py}.

The processing order, matching the code exactly, is shown in
Figure~\ref{fig:pod_pipeline_flow}.
Steps~2--3 are invertible and are undone during lifting (reconstruction).

\begin{figure}[h]
\centering
\begin{tikzpicture}[
    node distance=0.55cm and 0.0cm,
    stage/.style={
      draw, rounded corners=3pt, minimum width=11cm,
      minimum height=0.72cm, font=\small, align=center,
      fill=blue!6
    },
    arr/.style={-{Stealth[length=4pt]}, thick}
  ]
  \node[stage]                  (s1) {\textbf{1.}~Vectorise \& stack density movies $\;\to\;$ snapshot matrix $\mathbf{X}_{\mathrm{raw}}\in\mathbb{R}^{T_{\mathrm{tot}}\times d}$};
  \node[stage, below=of s1]    (s2) {\textbf{2.}~Shift-align all frames via FFT cross-correlation $\;\to\;$ $\mathbf{X}_{\mathrm{aligned}}$};
  \node[stage, below=of s2]    (s3) {\textbf{3.}~Point-wise density transform $\;\to\;$ $\tilde{\mathbf{X}} = \sqrt{\mathbf{X}_{\mathrm{aligned}}+\varepsilon}$};
  \node[stage, below=of s3]    (s4) {\textbf{4.}~Compute global mean $\boldsymbol{\mu}$; centre $\;\to\;$ $\bar{X} = \tilde{\mathbf{X}}^\top - \boldsymbol{\mu}\mathbf{1}^\top$};
  \node[stage, below=of s4]    (s5) {\textbf{5.}~Truncated SVD $\;\to\;$ POD basis $U_m\in\mathbb{R}^{d\times m}$};
  \node[stage, below=of s5]    (s6) {\textbf{6.}~Restrict: $\mathbf{y}_k = U_m^\top(\tilde{\mathbf{x}}_k - \boldsymbol{\mu})\;\in\mathbb{R}^{m}$};
  %
  \draw[arr] (s1) -- (s2);
  \draw[arr] (s2) -- (s3);
  \draw[arr] (s3) -- (s4);
  \draw[arr] (s4) -- (s5);
  \draw[arr] (s5) -- (s6);
\end{tikzpicture}
\caption{POD restriction pipeline: density fields to latent coordinates (six processing stages).}
\label{fig:pod_pipeline_flow}
\end{figure}

%
\section{Snapshot matrix construction}
\label{sec:snapshot_matrix}

We generate macroscopic density fields from microscopic particle trajectories
via KDE (Chapter~\ref{ch:background}, Algorithm~2) and assemble a global
snapshot matrix for POD.  Let $M$ denote the number of training runs and
$T_i$ the number of saved snapshots for run~$i$ (after temporal subsampling
by stride~$s$; \cv{subsample}${}=3$).

\subsection{Vectorisation and global stacking}
\label{subsec:vectorization}

Each density snapshot $\rho^{(i)}(t_k)\in\mathbb{R}^{n_y\times n_x}$ is
flattened into a column vector
\[
  \mathbf{s}^{(i)}_k := \mathrm{vec}\!\bigl(\rho^{(i)}(t_k)\bigr)
  \in\mathbb{R}^{d},\qquad d:=n_x\,n_y,
\]
and all frames from run~$i$ are stacked row-wise into
$\mathbf{X}^{(i)}\in\mathbb{R}^{T_i\times d}$.  The global training matrix
is the vertical concatenation
\[
  \mathbf{X}_{\mathrm{raw}}
  :=\begin{bmatrix}
      \mathbf{X}^{(0)}\\\vdots\\\mathbf{X}^{(M-1)}
    \end{bmatrix}
  \in\mathbb{R}^{T_{\mathrm{tot}}\times d},\qquad
  T_{\mathrm{tot}}=\sum_{i=0}^{M-1}T_i.
\]

\paragraph{Column-stacked notation.}
\label{subsec:column_notation}
Alvarez et~al.\ define a snapshot matrix whose \emph{columns} are density
snapshots: $X=[x_1,\dots,x_{n_t}]\in\mathbb{R}^{n_c\times n_t}$ with
$n_c{=}d$, $n_t{=}T_{\mathrm{tot}}$, and centre using $\bar{X}=XH$ where
$H=I-\tfrac{1}{n_t}\mathbf{1}\mathbf{1}^\top$~\cite{alvarez2025eqfree}.
To map our row-stacked implementation to this notation, set
$X:=\mathbf{X}_{\mathrm{raw}}^\top$.

\paragraph{Scope.}
The POD pipeline operates exclusively on the scalar density field~$\rho$.
The auxiliary fields used by WSINDy (flux $p_x,p_y$ and Morse potential~$\Phi$;
see Chapter~\ref{ch:background}, \S\ref{sec:multifield_cg}) are computed
in physical space and are \emph{not} projected through POD.

%
\section{Preprocessing: shift alignment and density transform}
\label{sec:preprocessing}

Two point-wise preprocessing steps are applied \emph{before} POD to improve
the compressibility of the snapshot manifold.

\subsection{Shift alignment}
\label{subsec:shift_align}

Travelling density patterns (flocks, mills) cause the dominant POD mode to
encode rigid translation rather than shape variation.  Shift alignment
removes this transport by registering every frame to a common reference via
circular cross-correlation.

\paragraph{Reference field.}
The reference is the temporal mean of all (unaligned) training snapshots,
reshaped to $n_y\times n_x$ (\cv{shift\_align\_ref: mean}):
\[
  \rho_{\mathrm{ref}}
  := \frac{1}{T_{\mathrm{tot}}}
     \sum_{k=1}^{T_{\mathrm{tot}}}
     \rho_k \;\in\mathbb{R}^{n_y\times n_x}.
\]

\paragraph{Per-frame shift estimation.}
For each snapshot $\rho_k$, the integer shift
$(\Delta y_k,\Delta x_k)$ is found by maximising the periodic
cross-correlation with the reference:
\begin{equation}
  (\Delta y_k,\Delta x_k)
  = \arg\max_{(j,i)}\;
    \mathcal{F}^{-1}\!\bigl[
      \mathcal{F}[\rho_{\mathrm{ref}}]
      \odot
      \overline{\mathcal{F}[\rho_k]}
    \bigr](j,i),
  \label{eq:shift_xcorr}
\end{equation}
where $\mathcal{F}$ denotes the 2-D DFT and the overline denotes complex
conjugation.  This is implemented in
\cv{shift\_align.fft\_cross\_correlation\_shift}.

\paragraph{Applying shifts.}
Each frame is circularly rolled by $(\Delta y_k,\Delta x_k)$ using
\cv{np.roll}, which is exact under periodic boundary conditions.  The shift
vectors are stored so that the inverse operation ($-\Delta y_k,-\Delta x_k$)
can undo the alignment during lifting.

\paragraph{Group-theoretic interpretation.}
The periodic domain $\Omega=[0,L_x)\times[0,L_y)$ is topologically a
\emph{flat torus}
$\mathbb{T}^2=(\mathbb{R}/L_x\mathbb{Z})\times(\mathbb{R}/L_y\mathbb{Z})$.
On the discrete grid, the cyclic translation group
$G=\mathbb{Z}_{n_y}\times\mathbb{Z}_{n_x}$ acts on density fields by
\[
  (\mathcal{T}_{\boldsymbol{\delta}}\,\rho)(i,j)
  \;=\;
  \rho\!\bigl((i-\delta_y)\bmod n_y,\;(j-\delta_x)\bmod n_x\bigr),
  \qquad
  \boldsymbol{\delta}=(\delta_y,\delta_x)\in G,
\]
which is exactly the circular shift implemented by \cv{np.roll}.  Each
$\mathcal{T}_{\boldsymbol{\delta}}$ is unitary on
$\ell^2(\mathbb{T}^2_d)$, so the action preserves both total mass and
$L^2$ energy.  The \emph{orbit}
$\mathrm{Orb}_G(\rho)=\{\mathcal{T}_g\rho:g\in G\}$ collects all
translates of a density field; two snapshots that differ only by a rigid
shift lie on the same orbit.

Shift alignment can now be stated as finding the group element that best
matches each frame to the reference:
\[
  \boldsymbol{\delta}_k^*
  = \arg\max_{\boldsymbol{\delta}\in G}
    \bigl\langle \rho_{\mathrm{ref}},\;
    \mathcal{T}_{\boldsymbol{\delta}}\,\rho_k\bigr\rangle,
\]
which is precisely the FFT cross-correlation in
Equation~\eqref{eq:shift_xcorr}.  Applying
$\mathcal{T}_{\boldsymbol{\delta}_k^*}$ to every snapshot
projects the data manifold~$\mathcal{M}$ onto its quotient
$\mathcal{M}/G$, where orbits collapse to single representatives.  Because
unaligned translation inflates the Kolmogorov $n$-width
of~$\mathcal{M}$~\cite{peherstorfer2022kolmogorov},
quotienting by~$G$ yields a sharply faster singular-value decay and thus
requires fewer POD modes for a given reconstruction accuracy.

This discrete single-reference alignment is a special case of the
\emph{shifted proper orthogonal decomposition} (sPOD) introduced by
Reiss et~al.~\cite{reiss2018shifted_pod}, which decomposes
transport-dominated fields into $K$ co-moving frames with time-dependent
shifts $\mathcal{T}_k\,q_k(x,t)=q_k(x-\Delta_k(t),t)$ and promotes
low-rank structure via nuclear-norm
minimisation~\cite{krah2024robust_spod,black2020projection}.
More broadly, our alignment echoes the symmetry-quotient paradigm
studied by Sonday et~al.~\cite{sonday2011alignment}, who use diffusion
maps to separate symmetry degrees of freedom from intrinsic dynamics in
the presence of noise, and by Budanur and
Cvitanovi\'{c}~\cite{budanur2015slices,siminos2015periodic}, whose
\emph{method of slices} defines a cross-section intersecting every
group orbit exactly once---a continuous analogue of the discrete
template-matching step above.

\subsection{Density transform}
\label{subsec:density_transform}

After alignment, a point-wise variance-stabilising transform is applied to
the (positive) density values:
\begin{equation}
  \tilde{\rho} := \sqrt{\rho + \varepsilon},
  \qquad \varepsilon = 10^{-10},
  \label{eq:sqrt_transform}
\end{equation}
(\cv{density\_transform: sqrt}, \cv{density\_transform\_eps}).  The square
root compresses the dynamic range of near-singular density spikes (e.g.\
NDYN05~blackhole) and improves the spectral decay of the snapshot
covariance, so that fewer POD modes are needed to capture a given energy
fraction.

All subsequent steps---mean computation, centering, SVD, restriction, and
lifting---operate in the \emph{transformed} ($\tilde\rho$) space.  The
inverse transform is applied after lifting:
\begin{equation}
  \hat{\rho} = \bigl[\max(0,\,\hat{\tilde\rho})\bigr]^2 - \varepsilon,
  \label{eq:sqrt_inverse}
\end{equation}
where the clamp to zero guards against small negative values introduced by
POD truncation.

\paragraph{Empirical role in the ROM comparison.}
In the preprocessing ablations, the density transform had little to no
effect on MVAR relative to working directly with raw density snapshots, but
it improved LSTM rollout stability and time-resolved $R^2$ more
consistently across regimes.  Since the thesis compares the mechanistic
WSINDy pipeline against black-box EF-ROM benchmarks, and compares MVAR and
LSTM on a shared latent representation, we retain the same density
transform for both surrogates for like-for-like benchmarking even though
the main benefit is observed on the LSTM side.

%
\section{POD via truncated SVD}
\label{sec:pod_svd}

Let $\mathbf{X}\in\mathbb{R}^{T_{\mathrm{tot}}\times d}$ denote the
shift-aligned, $\sqrt{\cdot}$-transformed snapshot matrix.  Compute the
global mean $\boldsymbol{\mu}=\tfrac{1}{T_{\mathrm{tot}}}
\mathbf{X}^\top\mathbf{1}\in\mathbb{R}^d$ and the centred matrix
\[
  \bar{X} := \mathbf{X}^\top - \boldsymbol{\mu}\,\mathbf{1}^\top
  \;\in\mathbb{R}^{d\times T_{\mathrm{tot}}}.
\]
Compute the thin SVD
\[
  \bar{X} = U\,\Sigma\,V^\top,
\]
where $U\in\mathbb{R}^{d\times r}$ contains the spatial POD modes,
$\Sigma=\mathrm{diag}(\sigma_1,\dots,\sigma_r)$, and $r=\mathrm{rank}(\bar{X})$.
This matches the construction in Alvarez et~al.~\cite{alvarez2025eqfree}

\paragraph{Energy spectrum and truncation.}
\label{subsec:energy_truncation}
Define the total energy $E_{\mathrm{tot}}=\sum_{k=1}^r\sigma_k^2$ and
cumulative explained ratio
\[
  \tau(m) := \frac{\sum_{k=1}^{m}\sigma_k^2}{E_{\mathrm{tot}}},
  \qquad 1\le m\le r.
\]
In production we use a \emph{fixed} latent dimension $m=19$
(\cv{fixed\_modes: 19}) and retain $U_m\in\mathbb{R}^{d\times m}$.

%
\section{Restriction operator \texorpdfstring{$\mathcal{R}$}{R}:
  density \texorpdfstring{$\to$}{to} latent coordinates}
\label{sec:restriction}

In the equation-free setting, restriction is the map
$\mathcal{R}:\mathbb{R}^d\to\mathbb{R}^m$ that takes a (transformed,
centred) snapshot to its latent
coordinates~\cite{alvarez2025eqfree}:
\begin{equation}
  \mathbf{y} = \mathcal{R}(\mathbf{x})
  := U_m^\top\bigl(\mathbf{x}-\boldsymbol{\mu}\bigr)
  \;\in\mathbb{R}^m,
  \label{eq:restriction}
\end{equation}
where $\mathbf{x}=\mathrm{vec}(\tilde\rho)$ is the transformed density.
In batch form,
$Y = U_m^\top(X - \boldsymbol{\mu}\,\mathbf{1}^\top)
\in\mathbb{R}^{m\times n_t}$.

In the row-stacked implementation (\cv{pod\_builder.py}), the equivalent
operation is
\[
  \mathbf{Y}_{\mathrm{row}} = \bigl(\mathbf{X}-\mathbf{1}\,\boldsymbol{\mu}^\top\bigr)\,U_m
  \;\in\mathbb{R}^{T_{\mathrm{tot}}\times m},
\]
so that each row of $\mathbf{Y}_{\mathrm{row}}$ is a latent snapshot.

%
\section{Lifting operator \texorpdfstring{$\mathcal{L}$}{L}:
  latent \texorpdfstring{$\to$}{to} reconstructed density}
\label{sec:lifting}

Lifting inverts the restriction by projecting back to the POD subspace and
adding the mean:
\begin{equation}
  \hat{\mathbf{x}} = \mathcal{L}(\mathbf{y})
  := U_m\,\mathbf{y} + \boldsymbol{\mu}.
  \label{eq:lifting}
\end{equation}
{\setlength{\abovedisplayskip}{5pt}
\setlength{\belowdisplayskip}{5pt}
\setlength{\abovedisplayshortskip}{4pt}
\setlength{\belowdisplayshortskip}{4pt}
\setlength{\jot}{3pt}
Since $\boldsymbol{\mu}$ and $U_m$ live in the transformed
($\sqrt{\cdot}$) space, the full reconstruction pipeline is:
\begin{enumerate}[leftmargin=*, itemsep=1pt, topsep=3pt, parsep=0pt, partopsep=0pt]
  \item Lift: $\hat{\mathbf{x}} = U_m\,\mathbf{y}+\boldsymbol{\mu}$
        \hfill(in $\sqrt{\rho}$ space)
  \item Inverse transform:
        $\hat\rho = [\max(0,\hat{\mathbf{x}})]^2 - \varepsilon$
        \hfill(Eq.~\ref{eq:sqrt_inverse})
  \item Optional clamp:
        $\hat\rho^{+} \leftarrow \max(\hat\rho,0)$
        \hfill(\cv{clamp\_mode: C1})
  \item Optional mass post-processing:
        $\hat\rho^\star \leftarrow \mathcal{P}_{M_0}(\hat\rho^{+})$
        or $\hat\rho^\star \leftarrow \mathcal{P}_{M_0}(\hat\rho)$
        \hfill(\cv{mass\_postprocess})
  \item Undo shift: $\hat\rho_{2\mathrm{D}}
        \leftarrow \mathrm{roll}(\hat\rho_{2\mathrm{D}},\;
        -\Delta y_k,-\Delta x_k)$
\end{enumerate}
Here \(M_0 := \sum_{j=1}^{d}\rho_j(t_{n_0})\) denotes the true total mass at
forecast start. Define
\begin{align}
  \Delta^{d-1}
  &:= \left\{
    p\in\mathbb{R}^{d} \,:\,
    p_j\ge 0,\;
    \sum_{j=1}^{d} p_j = 1
  \right\},
  \label{eq:probability_simplex}\\
  \Delta^{d-1}_{M_0}
  &:= M_0\,\Delta^{d-1}
  =
  \left\{
    \rho\in\mathbb{R}^{d} \,:\,
    \rho_j\ge 0,\;
    \sum_{j=1}^{d} \rho_j = M_0
  \right\}.
  \label{eq:mass_simplex}
\end{align}
Then the two mass-postprocessing operators used in the code are
\begin{align}
  \Pi_{\Delta^{d-1}_{M_0}}(\hat\rho)
  &:= \arg\min_{\rho\in\Delta^{d-1}_{M_0}}
  \|\rho-\hat\rho\|_2^2,
  \label{eq:mass_postprocess_simplex}\\
  \mathcal{S}_{M_0}(\hat\rho)
  &:=
  \begin{cases}
    \hat\rho\,\dfrac{M_0}{\sum_{j=1}^{d}\hat\rho_j},
    & \text{if }\sum_{j=1}^{d}\hat\rho_j > 0,\\[6pt]
    \hat\rho,
    & \text{if }\sum_{j=1}^{d}\hat\rho_j \le 0,
  \end{cases}
  \label{eq:mass_postprocess_scale}
\end{align}
}
which matches the implementation in \cv{test\_evaluator.py}: the frame is
rescaled only when its current total mass is positive.  This operator does
not itself enforce nonnegativity.  When \cv{scale} is paired with the
clamp-only correction \cv{C1}, the code first forms
$\hat\rho^{+}=\max(\hat\rho,0)$ and then applies
$\mathcal{S}_{M_0}(\hat\rho^{+})$; when no clamping is used, it acts directly
on $\hat\rho$.  By contrast, \cv{simplex} enforces both nonnegativity and the
mass constraint in a single step.  In the current preprocessing ablations,
either choice had little to no effect on MVAR forecast quality, but both
helped the LSTM maintain higher rollout $R^2$ and better mass control.  For
comparability, the same \cv{mass\_postprocess} choice is applied to both
black-box surrogates in any given experiment.  The final thesis results will
retain whichever of \cv{scale} or \cv{simplex} is chosen after the last
ablation round.  The final shift undoing in Step~5 restores the original
spatial position of the density pattern.

%
\section{Mass preservation under linear POD}
\label{sec:mass_preservation}

When no density transform is applied (i.e.\ the pipeline is purely linear),
POD lifting preserves total mass exactly.  We restate this result from
Alvarez et~al.~\cite{alvarez2025eqfree} in our notation.

\begin{proposition}[Mass preservation under constant-mass snapshots]
\label{prop:mass_preservation}
Let $X=[x_1,\dots,x_{n_t}]\in\mathbb{R}^{d\times n_t}$ with
$\mathbf{1}^\top x_k = M_0$ for all~$k$.
Let $\bar{X}=XH$, $\bar{X}=U\Sigma V^\top$, and define the reconstruction
$\tilde{X}:=U_m U_m^\top\bar{X}+\boldsymbol{\mu}\mathbf{1}^\top$.
Then $\mathbf{1}^\top\tilde{x}_k=M_0$ for all~$k$.
\end{proposition}

\begin{proof}
Constant mass implies $\mathbf{1}^\top X = M_0\mathbf{1}^\top$.  Hence
$\mathbf{1}^\top\bar{X}=\mathbf{1}^\top XH = M_0\mathbf{1}^\top H
=\mathbf{0}^\top$ because $H\mathbf{1}=\mathbf{0}$.
So $\mathbf{1}$ is orthogonal to every column of~$\bar{X}$ and therefore to
every left singular vector: $\mathbf{1}^\top U_m=\mathbf{0}^\top$.
Then
$\mathbf{1}^\top\tilde{X}
= \underbrace{\mathbf{1}^\top U_m}_{\mathbf{0}^\top} U_m^\top\bar{X}
  +\mathbf{1}^\top\boldsymbol{\mu}\,\mathbf{1}^\top
= M_0\mathbf{1}^\top$.
\end{proof}

\begin{remark}[Effect of the density transform]
The $\sqrt{\cdot}$ transform is nonlinear: even if
$\sum_j\rho_j=M_0$ for all frames, the transformed snapshots
$\tilde\rho_j=\sqrt{\rho_j+\varepsilon}$ do \emph{not} share a common mass.
Proposition~\ref{prop:mass_preservation} therefore does not apply in the
transformed space, and an optional mass post-processing operator
(\cv{mass\_postprocess: scale} or \cv{simplex}) is used to restore physical
mass after inverse transformation.
\end{remark}

%
\begin{algorithm}[H]
\caption{Global POD restriction--lifting pipeline (production).}
\label{alg:pod_pipeline}
\begin{algorithmic}[1]
\Require Training density movies $\{\rho^{(i)}(t_k)\}_{i=1}^M$, grid $(n_x,n_y)$,
  latent dim $m$, subsample stride $s$
\Ensure POD basis $U_m$, mean $\boldsymbol{\mu}$, shift data, latent trajectories
  $\{y^{(i)}(t_k)\}$
\Statex
\Statex \textbf{--- Training (restriction) ---}
\For{each run $i=1,\dots,M$}
  \State Subsample: keep every $s$-th frame
  \State Flatten: $x^{(i)}_k = \mathrm{vec}(\rho^{(i)}(t_k))\in\mathbb{R}^d$
\EndFor
\State Stack all runs: $\mathbf{X}_{\mathrm{raw}}\in\mathbb{R}^{T_{\mathrm{tot}}\times d}$
  \Comment{\cv{np.vstack}}
\State Compute reference $\rho_{\mathrm{ref}}$ (temporal mean of all frames)
  \Comment{\cv{shift\_align}}
\State Estimate per-frame shifts $(\Delta y_k,\Delta x_k)$ via FFT
  cross-correlation (Eq.~\ref{eq:shift_xcorr})
\State Apply shifts: $\rho_k\leftarrow\mathrm{roll}(\rho_k,\Delta y_k,\Delta x_k)$
\State Transform: $\tilde\rho_k\leftarrow\sqrt{\rho_k+\varepsilon}$
  \Comment{\cv{density\_transform: sqrt}}
\State Compute mean: $\boldsymbol{\mu}=\tfrac{1}{T_{\mathrm{tot}}}\sum_k \tilde\rho_k$
\State Centre: $\bar{X}=\mathbf{X}^\top-\boldsymbol{\mu}\mathbf{1}^\top$
\State SVD: $\bar{X}=U\Sigma V^\top$; keep $U_m$ ($m$ columns)
  \Comment{\cv{np.linalg.svd}}
\State Restrict: $Y=U_m^\top(\mathbf{X}^\top-\boldsymbol{\mu}\mathbf{1}^\top)$
\State Reshape $Y$ into per-run latent trajectories $\{y^{(i)}(t_k)\}$
\Statex
\Statex \textbf{--- Evaluation (lifting) ---}
\State Forecast latent trajectory $\hat{y}(t_k)$ (MVAR or LSTM; Chapter~\ref{ch:latent_rom})
\State Lift: $\hat{\tilde\rho}_k = U_m\,\hat{y}(t_k)+\boldsymbol{\mu}$
\State Inverse transform: $\hat\rho_k = [\max(0,\hat{\tilde\rho}_k)]^2-\varepsilon$
\State Optional mass postprocess: apply \eqref{eq:mass_postprocess_scale}
  or \eqref{eq:mass_postprocess_simplex}
  \Comment{\cv{mass\_postprocess: scale | simplex}}
\State Undo shift: $\hat\rho_k\leftarrow\mathrm{roll}(\hat\rho_k,-\Delta y_k,-\Delta x_k)$
\end{algorithmic}
\end{algorithm}
 
%
\chapter{Latent Reduced-Order Model: MVAR and LSTMs in POD Coordinates}
\label{ch:latent_rom}

This chapter describes the two latent-space time-steppers used in the
restrict--evolve--lift pipeline of Chapter~\ref{ch:global_pod_pipeline}:
a multivariate autoregressive model~(MVAR) and a long short-term memory
network~(LSTM).
Both models operate on the $r$-dimensional POD coefficient vector
$y(t_k)\in\mathbb{R}^r$ and produce a one-step-ahead prediction that is
iterated in closed loop to generate long-horizon density forecasts.

%
\section{MVAR(\texorpdfstring{$w$}{w}) model}
\label{sec:mvar_definition}
%

Let $x(t)$ denote the macroscopic density field on a fixed grid, sampled at
discrete times $t_k := t_0 + k\,\Delta t$ ($k = 0,1,2,\dots$) with effective sampling interval
$\Delta t = \texttt{dt}\times\texttt{subsample} = 0.04\times 3 = 0.12$\,s
(config keys \cv{dt} and \cv{subsample}).
The restriction operator $\mathcal{R}$ (Chapter~\ref{ch:global_pod_pipeline}) maps
$x(t_k)$ to a latent coordinate $y(t_k)\in\mathbb{R}^r$ with
$r=19$ (\cv{fixed\_modes}).
In this chapter we focus on learning a surrogate time-stepper for $y$.

\subsection{Model class}
We model the latent coordinates as a discrete-time vector autoregression with
$w$ lags (\cv{lag}).
For each training trajectory (``case'') $c$ we write
\begin{equation}\label{eq:mvar_def}
  y^{(c)}(t_k) \;=\; A_0 \;+\; \sum_{j=1}^{w} A_j\,y^{(c)}(t_{k-j})
                 \;+\; \varepsilon(t_k),
\end{equation}
where $A_0\in\mathbb{R}^r$ is an intercept vector and
$A_j\in\mathbb{R}^{r\times r}$ are lag-$j$ regression matrices~\cite{alvarez2025eqfree}.
The residual $\varepsilon(t_k)\in\mathbb{R}^r$ is modelled as zero-mean
white noise; this assumption motivates least-squares training but is not
used for inference (we do point prediction).

\paragraph{One-step predictor.}
The corresponding one-step-ahead prediction is
\begin{equation}\label{eq:mvar_predictor}
  \hat y^{(c)}(t_k) \;=\; A_0 \;+\;
  \sum_{j=1}^{w} A_j\,y^{(c)}(t_{k-j}).
\end{equation}

\subsection{Ridge-regularized training}\label{subsec:mvar_training}
The unknown parameters $(A_0,\{A_j\}_{j=1}^w)$ are obtained by minimizing
\begin{equation}\label{eq:mvar_ls_problem}
  \min_{A_0,\{A_j\}_{j=1}^w}\;
  \sum_{c=1}^{C}\;\sum_{k=w+1}^{K}
  \bigl\|y^{(c)}(t_k)-\hat y^{(c)}(t_k)\bigr\|_2^2.
\end{equation}

\paragraph{Vectorized regression form.}
Define the stacked regressor
\[
  \phi^{(c)}_k \;:=\;
  \begin{bmatrix}1\\[2pt] y^{(c)}(t_{k-1})\\ \vdots\\ y^{(c)}(t_{k-w})\end{bmatrix}
  \!\in\mathbb{R}^{1+wr},
\]
and let
$B\in\mathbb{R}^{(1+wr)\times r}$ collect the intercept and lag matrices
so that $\hat y^{(c)}(t_k) = (\phi^{(c)}_k)^\top B$.
Stacking all $N_{\mathrm{train}}=C(K-w)$ training pairs into matrices
$X\in\mathbb{R}^{N_{\mathrm{train}}\times(1+wr)}$ and
$Y\in\mathbb{R}^{N_{\mathrm{train}}\times r}$,
ridge-regularized least squares solves
\begin{equation}\label{eq:mvar_ridge}
  \hat B \;=\; \arg\min_{B}\;\|Y - XB\|_F^2
             \;+\;\alpha\,\|B_{2:}\|_F^2,
\end{equation}
where $B_{2:}$ excludes the intercept row from regularization and
$\alpha>0$ (\cv{ridge\_alpha}).
The closed-form solution is
\begin{equation}\label{eq:ridge_closed_form}
  \hat B \;=\; (X^\top X + \alpha\,\Gamma)^{-1}\,X^\top Y,
  \qquad
  \Gamma = \mathrm{diag}(0,1,\dots,1).
\end{equation}
In our codebase, this is implemented in
\cv{train\_mvar\_model()} \mbox{(\cv{src/rectsim/mvar\_trainer.py})}
using \texttt{sklearn.linear\_model.Ridge} with
\texttt{fit\_intercept=True}, which automatically excludes the intercept
from the penalty term, matching~\eqref{eq:ridge_closed_form}.

\paragraph{Production values.}
In the baseline systematic comparison, the MVAR uses lag
$w_{\mathrm{MVAR}}=5$ (\cv{lag:~5}) and ridge strength
$\alpha=10^{-4}$ (\cv{ridge\_alpha:~0.0001}).  At
$\Delta t_{\mathrm{rom}}=0.12$\,s this lag window spans $0.6$\,s of
physical time.  This is a compromise between capturing short-term latent
dependence and keeping the regression well-conditioned; ridge
regularization provides additional robustness.  When the LSTM lag is
varied in dedicated ablations, the evaluation protocol uses a common
forecast start determined by the largest enabled lag in the config.
Systematic lag selection via AIC/BIC~\cite{alvarez2025eqfree} is left as
future work.

\subsection{Closed-loop rollout}\label{subsec:mvar_rollout}
For long-horizon forecasting the MVAR model is used in \emph{closed-loop}
(autoregressive) mode: each one-step prediction is fed back as input for
subsequent steps.

\paragraph{Rollout map.}
Let the state of the autoregression be
\[
  z_k :=
  \begin{bmatrix}
    y(t_k)\\ y(t_{k-1})\\ \vdots\\ y(t_{k-w+1})
  \end{bmatrix}\in\mathbb{R}^{wr}.
\]
The MVAR($w$) model induces an affine update
\begin{equation}\label{eq:companion_update}
  z_{k+1} = \mathcal{A}_{\mathrm{comp}}\, z_k + b_{\mathrm{comp}},
\end{equation}
where $\mathcal{A}_{\mathrm{comp}}\in\mathbb{R}^{wr\times wr}$ is the
companion matrix: the top block row contains $(A_1,\dots,A_w)$ and the
subdiagonal contains identity ``shift'' blocks.
Errors compound through recursion, so closed-loop metrics are the primary
evaluation tool (Chapter~\ref{ch:exp_protocol_metrics}).

\paragraph{Restrict--evolve--lift forecast.}
In the full pipeline the autoregressive forecast operates on physical-space
densities:
\begin{equation}\label{eq:closed_loop_forecast}
  \hat x^{(c)}(t_k) \;=\;
  \mathcal{L}\!\Bigl(
    \Phi_{\mathrm{ROM}}\!\bigl(
      \mathcal{R}(\hat x^{(c)}(t_{k-1})),\dots,
      \mathcal{R}(\hat x^{(c)}(t_{k-w}))
    \bigr)
  \Bigr),
\end{equation}
where $\mathcal{R}$ and $\mathcal{L}$ are the restriction and lifting
operators of Chapter~\ref{ch:global_pod_pipeline}, and $\Phi_{\mathrm{ROM}}$ is
either the MVAR or LSTM time-stepper.
The forecast is initialized with the first $w$ true (restricted) snapshots
and iterated for all remaining time steps.

\paragraph{Stability.}
Closed-loop performance depends on the spectral radius
\[
  \rho(\mathcal{A}_{\mathrm{comp}}) < 1,
\]
which is a sufficient condition for bounded rollout of the homogeneous part.
Ridge regularization shrinks coefficients and therefore tends to improve
stability margins.

\paragraph{Optional spectral-radius scaling.}
The codebase includes a post-fit stabilization step controlled by
\cv{eigenvalue\_threshold}: if
$\rho(\mathcal{A}_{\mathrm{comp}})$ exceeds the threshold, all coefficient
matrices are uniformly scaled down.
In preliminary experiments this scaling degraded forecast quality by
over-damping transient dynamics, so it is \emph{disabled} in all production
runs.
Ridge regularization alone provides sufficient implicit stability control.

\paragraph{Parameter count.}\label{subsec:mvar_identifiability}
The total number of free parameters is
\[
  p \;=\; r + w\,r^2.
\]
With $r=19$ and $w=5$ this gives $p = 19 + 5\times 19^2 = 1824$
parameters, fitted from $N_{\mathrm{train}} = C(K{-}w)$ samples across all
training trajectories.
Fixing a modest POD dimension and using ridge regularization keeps the
effective parameter-to-data ratio manageable.

%
\section{LSTM latent surrogate}
\label{sec:lstm_surrogate}
%

The LSTM serves as the \emph{nonlinear} latent time-stepper in the EF-ROM
benchmark.
It provides a high-capacity predictor against which we compare the
interpretable WSINDy models.
The implementation lives in the class
\cv{LatentLSTMROM} (\cv{src/rom/lstm\_rom.py}), trained by the function
\cv{train\_lstm\_rom()}.

\paragraph{Sequence-to-one formulation.}
Let $\mathbf{z}_k\in\mathbb{R}^d$ denote the POD coefficient vector at
time $t_k$ ($d = r = 19$).
A $p$-step autoregressive surrogate learns the one-step map
\begin{equation}\label{eq:lstm_seq_problem}
  \mathbf{z}_{k} \;\approx\;
  f_\theta\!\bigl(\mathbf{z}_{k-1},\ldots,\mathbf{z}_{k-p}\bigr),
  \qquad k = p{+}1,\ldots,K,
\end{equation}
where $f_\theta$ is a parameterized nonlinear model.
LSTMs are preferred over vanilla RNNs because the gated cell state
mitigates vanishing/exploding gradients, which matters when the latent
series exhibits long transients or regime shifts.

\paragraph{LSTM gate equations.}
Given an input window
$\mathbf{x}_k = (\mathbf{z}_{k-p},\ldots,\mathbf{z}_{k-1})
\in\mathbb{R}^{p\times d}$,
the network processes the sequence recursively:
\begin{align}
\mathbf{f}_\tau &= \sigma\!\bigl(\mathbf{W}_f[\mathbf{h}_{\tau-1},\mathbf{z}_\tau]+\mathbf{b}_f\bigr), &
\mathbf{i}_\tau &= \sigma\!\bigl(\mathbf{W}_i[\mathbf{h}_{\tau-1},\mathbf{z}_\tau]+\mathbf{b}_i\bigr), \nonumber\\
\mathbf{o}_\tau &= \sigma\!\bigl(\mathbf{W}_o[\mathbf{h}_{\tau-1},\mathbf{z}_\tau]+\mathbf{b}_o\bigr), &
\tilde{\mathbf{c}}_\tau &= \tanh\!\bigl(\mathbf{W}_c[\mathbf{h}_{\tau-1},\mathbf{z}_\tau]+\mathbf{b}_c\bigr), \nonumber\\
\mathbf{c}_\tau &= \mathbf{f}_\tau \odot \mathbf{c}_{\tau-1} + \mathbf{i}_\tau \odot \tilde{\mathbf{c}}_\tau, &
\mathbf{h}_\tau &= \mathbf{o}_\tau \odot \tanh(\mathbf{c}_\tau),
\label{eq:lstm_gates}
\end{align}
where $(\mathbf{h}_\tau,\mathbf{c}_\tau)$ are the hidden and cell states.
Figure~\ref{fig:lstm_cell} illustrates the information flow inside a single
LSTM cell.

\begin{figure}[h]
\centering
\begin{tikzpicture}[
    >=Stealth, thick,
    gate/.style={draw, circle, minimum size=0.8cm, inner sep=0pt,
                 font=\footnotesize\bfseries, fill=orange!15},
    op/.style={draw, circle, minimum size=0.6cm, inner sep=0pt,
               font=\footnotesize, fill=blue!10},
    act/.style={draw, rectangle, rounded corners=2pt,
                minimum width=1.0cm, minimum height=0.55cm,
                font=\footnotesize, fill=green!10},
    io/.style={font=\small},
    lbl/.style={font=\scriptsize, text=black!60},
    line/.style={thick},
  ]

  %
  \node[io] (cin)   at (-0.5,3.6) {$\mathbf{c}_{\tau-1}$};
  \node[op] (mulf)  at (2.2,3.6)  {$\odot$};
  \node[op] (plus)  at (5.0,3.6)  {$+$};
  \node[io] (cout)  at (10.5,3.6) {$\mathbf{c}_{\tau}$};

  \draw[line,->] (cin)  -- (mulf);
  \draw[line,->] (mulf) -- (plus);
  \draw[line,->] (plus) -- (cout);

  %
  \node[gate] (fg)   at (2.2,1.8)  {$\sigma$};
  \node[gate] (ig)   at (4.0,1.8)  {$\sigma$};
  \node[act]  (cand) at (5.8,1.8)  {tanh};
  \node[gate] (og)   at (7.8,1.8)  {$\sigma$};

  %
  \node[lbl, left=1pt of fg]    {$\mathbf{f}$};
  \node[lbl, left=1pt of ig]    {$\mathbf{i}$};
  \node[lbl, right=1pt of cand] {$\tilde{\mathbf{c}}$};
  \node[lbl, left=1pt of og]    {$\mathbf{o}$};

  %
  \node[op] (muli) at (5.0,2.7) {$\odot$};

  %
  \coordinate (cbranch) at (7.8,3.6);
  \fill (cbranch) circle (2pt);            %
  \node[act] (thc)  at (7.8,2.7) {tanh};
  \node[op]  (mulo) at (9.2,1.8) {$\odot$};

  %
  \node[io] (hin)  at (-0.5,-0.2)  {$\mathbf{h}_{\tau-1}$};
  \node[io] (hout) at (10.5,-0.2)  {$\mathbf{h}_{\tau}$};
  \node[io] (zt)   at (4.8,-1.5)   {$\mathbf{z}_{\tau}$};

  %
  \node[act, fill=gray!10] (concat) at (4.8,0.4)
    {$[\,\cdot\,,\,\cdot\,]$};

  %
  \draw[line,->] (hin) |- (concat);
  \draw[line,->] (zt)  -- (concat);

  %
  \coordinate (fan) at (4.8,1.0);
  \draw[line]    (concat) -- (fan);
  \draw[line,->] (fan) -| (fg);
  \draw[line,->] (fan) -| (ig);
  \draw[line,->] (fan) -| (cand);
  \draw[line,->] (fan) -| (og);

  %
  %
  \draw[line,->] (fg) -- (mulf);

  %
  \draw[line,->] (ig)   |- (muli);
  \draw[line,->] (cand) |- (muli);
  \draw[line,->] (muli) -- (plus);

  %
  %
  \draw[line,->] (cbranch) -- (thc);
  \draw[line,->] (thc)  -| (mulo);
  \draw[line,->] (og)   -- (mulo);
  \draw[line,->] (mulo) |- (hout);

\end{tikzpicture}
\caption{Single LSTM cell.
The cell state~$\mathbf{c}$ flows along the top highway (left to right);
the \emph{forget} gate~($\mathbf{f}$) erases selectively via~$\odot$,
the \emph{input} gate~($\mathbf{i}$) and candidate~($\tilde{\mathbf{c}}$)
write new information via~$\odot$ then~$+$,
and the \emph{output} gate~($\mathbf{o}$) controls what enters the hidden
state~$\mathbf{h}_\tau$ through $\tanh(\mathbf{c}_\tau)\odot\mathbf{o}$.
Inputs $[\mathbf{h}_{\tau-1},\mathbf{z}_\tau]$ are concatenated and fed
to all four learnable transformations.}
\label{fig:lstm_cell}
\end{figure}

\paragraph{Layer normalisation and readout.}
The final hidden state $\mathbf{h}_{p}$ is passed through a
\texttt{LayerNorm} layer (\cv{use\_layer\_norm:~true}) before the linear
readout:
\begin{equation}\label{eq:lstm_readout}
  \widehat{\mathbf{z}}_{k} \;=\;
  \mathbf{W}_y\,\mathrm{LayerNorm}(\mathbf{h}_{p}) + \mathbf{b}_y.
\end{equation}

\paragraph{Residual (delta) prediction.}
In production, the network is configured with \cv{residual:~true}, so
the output layer predicts a \emph{correction} $\Delta_k$ and the final
prediction is
\begin{equation}\label{eq:lstm_residual}
  \widehat{\mathbf{z}}_{k} \;=\;
  \mathbf{z}_{k-1} \;+\; \Delta_k,
  \qquad
  \Delta_k := \mathbf{W}_y\,\mathrm{LayerNorm}(\mathbf{h}_{p}) + \mathbf{b}_y.
\end{equation}
This residual connection anchors the prediction to the most recent input
state, which improves stability during closed-loop rollout.

\paragraph{Training objective.}
The LSTM is trained via one-step prediction on sliding windows of
ground-truth latent states.  Each training sample is an independent
(window, target) pair; the LSTM processes the full $p$-step window in a
single forward pass and outputs one prediction---there is no sequential
decoding step where the model's own output is fed back, so the term
``teacher forcing'' does not apply.
The loss minimised is mean-squared error:
\begin{equation}\label{eq:lstm_mse}
  \min_{\theta}\;
  \frac{1}{N_{\mathrm{win}}}\sum_{\ell=1}^{N_{\mathrm{win}}}
  \bigl\|\widehat{\mathbf{z}}_{k(\ell)}-\mathbf{z}_{k(\ell)}\bigr\|_2^2,
\end{equation}
optimized with Adam.
Early stopping on a held-out validation split (\cv{patience:~30}),
gradient clipping (\cv{grad\_clip:~1.0}), and a
\texttt{ReduceLROnPlateau} scheduler (factor $0.5$, patience~$20$)
are used during training.

\paragraph{Production configuration.}
Table~\ref{tab:lstm_config} lists the LSTM hyperparameters used in all
experiments, as specified in the \cv{rom.models.lstm} block of each
configuration file.

\begin{table}[h]
\centering\small
\caption{LSTM hyperparameters (production configuration).}
\label{tab:lstm_config}
\begin{tabular}{lll}
\toprule
\textbf{Parameter} & \textbf{Config key} & \textbf{Value} \\
\midrule
Lag window          & \cv{lag}              & 5 (baseline) \\
Hidden units        & \cv{hidden\_units}    & 128 \\
LSTM layers         & \cv{num\_layers}      & 2   \\
Dropout (inter-layer) & \cv{dropout}        & 0.1 \\
Residual connection & \cv{residual}         & true \\
Layer normalisation & \cv{use\_layer\_norm} & true \\
Max epochs          & \cv{max\_epochs}      & 300 \\
Early-stop patience & \cv{patience}         & 30  \\
Learning rate       & \cv{lr}               & $10^{-3}$ \\
Batch size          & \cv{batch\_size}      & 256 \\
Gradient clip       & \cv{grad\_clip}       & 1.0 \\
\bottomrule
\end{tabular}
\end{table}

\paragraph{Closed-loop forecasting.}
After training, we forecast by recursively feeding predictions back into
the input window:
\[
  (\mathbf{z}_{k-p},\ldots,\mathbf{z}_{k-1})
  \;\mapsto\; \widehat{\mathbf{z}}_{k}
  \;\mapsto\; (\mathbf{z}_{k-p+1},\ldots,\mathbf{z}_{k-1},\widehat{\mathbf{z}}_{k})
  \;\mapsto\;\cdots
\]
This closed-loop evaluation matches the restrict--evolve--lift protocol of
\eqref{eq:closed_loop_forecast}.
Unlike MVAR, the LSTM has no eigenvalue-based stability certificate;
stability is monitored empirically via rollout diagnostics
(boundedness of POD coefficients and mass drift after lifting).

\paragraph{Algorithmic summary.}
Algorithm~\ref{alg:lstm_rom} summarises the LSTM ROM pipeline.

\begin{algorithm}[H]
\caption{Latent LSTM ROM (\cv{LatentLSTMROM}): training and closed-loop rollout}
\label{alg:lstm_rom}
\begin{algorithmic}[1]
\Require Latent trajectories $\{y^{(c)}_k\}$, lag $p$ (\cv{lag}), horizon $H$
\Ensure Closed-loop forecast $\{\widehat{\mathbf{z}}_{k_0+n}\}_{n=1}^{H}$

\State \textbf{Windowing:} construct pairs
  $(\mathbf{x}_k, \mathbf{z}_k)$ for $k=p{+}1,\ldots,K$,
  with $\mathbf{x}_k=(\mathbf{z}_{k-p},\ldots,\mathbf{z}_{k-1})$.
\State Split into train/validation; initialise LSTM weights.

\State \textbf{Training} (Adam, lr${}=10^{-3}$)\textbf{:}
\While{epoch $\le 300$ and no early stop (patience${}=30$)}
  \State Sample mini-batch of size $B=256$.
  \State Forward pass through \eqref{eq:lstm_gates}; apply LayerNorm to $\mathbf{h}_p$.
  \State Compute $\widehat{\mathbf{z}}_k = \mathbf{z}_{k-1} + \Delta_k$ \quad (residual mode, \eqref{eq:lstm_residual}).
  \State Backpropagate MSE loss \eqref{eq:lstm_mse}; clip gradients at $1.0$.
  \State Step \texttt{ReduceLROnPlateau} scheduler on validation loss.
\EndWhile

\State \textbf{Closed-loop rollout:}
\State Set window $\mathbf{x}\gets\mathbf{x}_{k_0}$.
\For{$n=1,\ldots,H$}
  \State Predict $\widehat{\mathbf{z}}_{k_0+n}$ from $\mathbf{x}$ (residual mode).
  \State Shift window; append $\widehat{\mathbf{z}}_{k_0+n}$.
\EndFor
\State \textbf{Lift} each $\widehat{\mathbf{z}}_{k_0+n}$ to physical space
  via the POD basis (Chapter~\ref{ch:global_pod_pipeline}).
\end{algorithmic}
\end{algorithm}

\paragraph{Comparison to MVAR.}
MVAR offers fast closed-form training and transparent stability analysis
(companion-matrix spectral radius), while LSTM provides nonlinear
capacity at the cost of more hyperparameters and weaker a~priori rollout
guarantees.
In crowd-dynamics EF-ROM benchmarks, linear MVARs can match or outperform
LSTMs in predictive accuracy while being lower-complexity and more
interpretable~\cite{alvarez2025eqfree}.
We include both in our pipeline: MVAR as a strong linear baseline, and
LSTM as a flexible nonlinear baseline relevant when latent dynamics exhibit
clear nonlinearities (e.g.\ variable-speed or alignment-coupling effects).
 
%
\chapter{ROM Results: MVAR vs LSTM}
\section{Data regime: snapshots, windows, effective samples}
\section{Compute/runtime comparison}
\section{Qualitative results: density snapshots (truth vs prediction)}
\section{Quantitative results: $R^2(t)$ curves by IC type}
\section{Order parameter trajectories (truth vs prediction)}
\section{Stability analysis outcomes (before/after rescaling)}
\section{Discussion: why MVAR can fail despite success in \cite{alvarez2025eqfree}}
\subsection{Linearity + stationarity mismatch}
\subsection{Heavy-tail POD spectrum and latent dimension burden}
\subsection{Data-to-parameter ratio and lag/window constraints}
\section{When LSTM improves and remaining failure modes}
\section{What worked, what did not}
\section{Immediate next steps before Part II}
 
%
\part{Weak-form PDE Discovery (WSINDy)}

%
%
%
%
\chapter{WSINDy: Weak-Form PDE Discovery}
\label{ch:wsindy_methodology}

This chapter presents the mathematical methodology of the
\emph{Weak-form Sparse Identification of Nonlinear Dynamics} (WSINDy)
pipeline~\cite{messenger2021wsindy} used to discover macroscopic PDE
models directly from the coarse-grained fields generated from the same
regime-specific simulation family and train/test split used by the EF-ROM
benchmarks in Part~I\@.  The data are common only upstream, however.
Whereas Part~I shifts, transforms, and compresses density snapshots into POD
coordinates and then learns a discrete-time latent evolution map, WSINDy stays
in physical field space and seeks an explicit continuous-time PDE of the form
\begin{equation}
  u_t = \sum_{m=1}^{M} w_m\,\mathcal{L}_m\!\bigl[f_m(u)\bigr],
  \label{eq:wsindy_general}
\end{equation}
where each candidate term is a (differential-operator, nonlinear-feature)
pair $(\mathcal{L}_m, f_m)$ drawn from a user-specified library, and the
sparse coefficient vector $\mathbf{w}\in\mathbb{R}^M$ is to be determined.
The pipeline is implemented in \cv{src/wsindy/} and orchestrated by
\cv{ROM\_WSINDY\_pipeline.py}.

The processing stages are shown in Figure~\ref{fig:wsindy_pipeline_flow}.
Sections~\ref{sec:weak_form}--\ref{sec:forecast} describe each stage
in order.
\ifdefined\feedbackdraft
For this feedback draft, only Sections~\ref{sec:weak_form}--\ref{sec:weak_system}
are filled in; Sections~\ref{sec:regression}--\ref{sec:forecast} are left as
placeholders for later revision.
\else
Detailed pseudocode for the two most intricate components—the sparse
regression algorithm and the time integrator—is deferred to
Appendix~\ref{app:mstls} and Appendix~\ref{app:etdrk4}, respectively.
\fi

\begin{figure}[ht]
\centering
\begin{tikzpicture}[
    node distance=0.55cm and 0.0cm,
    stage/.style={
      draw, rounded corners=3pt, minimum width=12cm,
      minimum height=0.72cm, font=\small, align=center,
      fill=blue!6
    },
    arr/.style={-{Stealth[length=4pt]}, thick}
  ]
  \node[stage]                  (s1) {\textbf{1.}~Construct separable test functions $\psi_k$ and derivative tensors (\S\ref{sec:weak_form})};
  \node[stage, below=of s1]    (s2) {\textbf{2.}~Build candidate library: operators $\times$ features (\S\ref{sec:library})};
  \node[stage, below=of s2]    (s3) {\textbf{3.}~Assemble weak linear system $\mathbf{b}=\mathbf{G}\mathbf{w}$ via FFT convolution (\S\ref{sec:weak_system})};
  \node[stage, below=of s3]    (s4) {\textbf{4.}~MSTLS sparse regression $+$ model selection over $\boldsymbol\ell$ grid (\S\ref{sec:regression})};
  \node[stage, below=of s4]    (s5) {\textbf{5.}~Bootstrap UQ $+$ stability selection (\S\ref{sec:uq})};
  \node[stage, below=of s5]    (s6) {\textbf{6.}~Forecast: ETDRK4 / RK4 integration of discovered PDE (\S\ref{sec:forecast})};
  \draw[arr] (s1) -- (s2);
  \draw[arr] (s2) -- (s3);
  \draw[arr] (s3) -- (s4);
  \draw[arr] (s4) -- (s5);
  \draw[arr] (s5) -- (s6);
\end{tikzpicture}
\caption{WSINDy pipeline: from density snapshots to a discovered PDE and its forecast.}
\label{fig:wsindy_pipeline_flow}
\end{figure}

%
\section{The weak-form framework}
\label{sec:weak_form}

Standard sparse-regression approaches (e.g.\ SINDy) evaluate spatial
derivatives of the data numerically, amplifying measurement noise.
WSINDy avoids this by reformulating the identification problem in
\emph{weak form}: derivatives are transferred from the noisy data~$u$ onto
smooth, compactly supported test functions via integration by parts.

\subsection{Test functions}
\label{subsec:test_functions}

We use separable polynomial bump functions
\begin{equation}
  \psi(x,y,t)
  = \phi_x(x)\;\phi_y(y)\;\phi_t(t),
  \label{eq:test_separable}
\end{equation}
where each factor is a compactly supported polynomial on
$[-\ell_i\Delta_i,\;\ell_i\Delta_i]$:
\begin{equation}
  \phi_i(s)
  = \begin{cases}
      \displaystyle
      \Bigl(1 - \frac{s^2}{(\ell_i\Delta_i)^2}\Bigr)^{\!p_i},
      & |s| \le \ell_i\Delta_i, \\[6pt]
      0, & \text{otherwise},
    \end{cases}
  \qquad i\in\{x,y,t\}.
  \label{eq:test_1d}
\end{equation}
Here $\ell_i$ is the half-width in grid cells, $\Delta_i$ the grid spacing
in direction~$i$, and $p_i$ the polynomial exponent (default $p_i=2$).
Increasing~$p_i$ yields smoother bumps at the cost of narrower effective
support; increasing~$\ell_i$ widens the averaging window.
In the implementation, the spacings are stored as
\cv{GridSpec(dt, dx, dy)}, so $(\Delta_t,\Delta_x,\Delta_y)$ correspond to
$(dt,dx,dy)$, and the integer half-widths are passed as
(\cv{ellt}, \cv{ellx}, \cv{elly}) to \cv{make\_separable\_psi}.
The implementation is in \cv{src/wsindy/test\_functions.py}
(\cv{make\_1d\_phi}).

Figure~\ref{fig:test_bump} shows the one-dimensional bump
$\phi(s)=(1-s^2)^p$ for several exponents.

\begin{figure}[ht]
\centering
\begin{tikzpicture}
\begin{axis}[
  width=10cm, height=5.5cm,
  domain=-1:1, samples=100,
  xlabel={$s/(\ell\Delta)$},
  ylabel={$\phi(s)$},
  ymin=0, ymax=1.12,
  legend style={at={(0.98,0.95)}, anchor=north east, font=\footnotesize},
  grid=major,
  grid style={gray!25},
  thick
]
\addplot[blue, thick]   {(1-x^2)^1};  \addlegendentry{$p=1$}
\addplot[red, thick]    {(1-x^2)^2};  \addlegendentry{$p=2$ (default)}
\addplot[orange, thick] {(1-x^2)^4};  \addlegendentry{$p=4$}
\end{axis}
\end{tikzpicture}
\caption{One-dimensional polynomial bump function
$\phi(s)=(1-s^2/(\ell\Delta)^2)^{p}$ normalised to unit half-width.
Higher~$p$ concentrates mass near the centre.}
\label{fig:test_bump}
\end{figure}

\subsection{Integration by parts}
\label{subsec:ibp}

Consider the PDE $u_t = F(u)$ on the periodic spatial domain
$\Omega=[0,L_x)\times[0,L_y)$ (\cv{Lx}, \cv{Ly}) over the time interval
$[0,T]$ (\cv{T}).  For a test function $\psi_k$ centred at a query point
$(x_k,y_k,t_k)$, let $Q_k := \operatorname{supp}\psi_k
\subset \Omega\times[0,T]$ denote its spacetime support.  In code,
$Q_k$ has half-widths $(\ell_t\Delta t,\ell_x\Delta x,\ell_y\Delta y)$,
corresponding to the grid variables \cv{dt}, \cv{dx}, \cv{dy} and the
integer half-widths \cv{ellt}, \cv{ellx}, \cv{elly}.  Multiplying both sides
by $\psi_k$ and integrating over~$Q_k$ gives
\[
  \int_{Q_k} \psi_k(x,y,t)\,u_t(x,y,t)\;dx\,dy\,dt
  = \int_{Q_k} \psi_k(x,y,t)\,F(u)(x,y,t)\;dx\,dy\,dt.
\]
Equivalently, since $\psi_k$ vanishes outside~$Q_k$, one may view these as
integrals over the full spacetime domain $\Omega\times[0,T]$.  Integrating
the left side by parts in time, the temporal boundary term vanishes because
$\psi_k$ is compactly supported in time:
\begin{equation}
  -\int_{Q_k} \frac{\partial\psi_k}{\partial t}(x,y,t)\,u(x,y,t)\;dx\,dy\,dt
  = \int_{Q_k} \psi_k(x,y,t)\,F(u)(x,y,t)\;dx\,dy\,dt.
  \label{eq:ibp_temporal}
\end{equation}
Likewise, for differential operators appearing in~$F(u)$, spatial
integration by parts transfers derivatives from the data~$u$ onto the test
function~$\psi_k$.  Under periodic boundary conditions on~$\Omega$, the
spatial boundary contributions vanish.  For example, if $F(u)$ contains a
term $w_m\,\Delta u$, then
\[
  w_m\int_{Q_k}\psi_k\,\Delta u\;dx\,dy\,dt
  = -\,w_m\int_{Q_k}\nabla\psi_k\cdot\nabla u\;dx\,dy\,dt
  = w_m\int_{Q_k}(\Delta\psi_k)\,u\;dx\,dy\,dt,
\]
where the intermediate equality follows from two spatial integrations by
parts with periodic boundary terms vanishing.  Thus the data~$u$ always appears
\emph{undifferentiated}.  The one-dimensional bump formula
\eqref{eq:test_1d} is analytic, but the sampled derivative arrays used in
the weak system are evaluated numerically on the discrete support lattice
via one-dimensional finite-difference stencils and then assembled by
separable outer products, as implemented in
\cv{make\_separable\_psi}.

\subsection{Separability and FFT convolution}
\label{subsec:fft_conv}

Because $\psi_k$ is separable~\eqref{eq:test_separable}, each
three-dimensional integral reduces to a product of one-dimensional
convolutions, computed efficiently via the Fast Fourier Transform.
The spatial dimensions are periodic and use circular convolution;
the temporal dimension is zero-padded to avoid wrap-around artefacts.
The implementation is in \cv{src/wsindy/system.py}
(\cv{fft\_convolve3d\_same}).

%
\section{Candidate library construction}
\label{sec:library}

The library specifies which physical mechanisms the regression is allowed
to select.  Each candidate term is a pair $(\mathcal{L}_m, f_m)$
consisting of a differential operator~$\mathcal{L}_m$ and a pointwise
nonlinear feature~$f_m$.

\subsection{Scalar library}
\label{subsec:scalar_library}

For a single field~$u$ the default library is the Cartesian product of
polynomial features and spatial operators:
\begin{equation}
  \bigl\{f_j\bigr\} = \{1,\;u,\;u^2,\;u^3\},
  \qquad
  \bigl\{\mathcal{L}_j\bigr\}
  = \{I,\;\partial_x,\;\partial_y,\;\Delta\},
  \label{eq:scalar_library}
\end{equation}
yielding up to 16 candidate terms (e.g.\ $\Delta u^2$, $\partial_x u$).
The maximum polynomial power is controlled by $p_{\max}$ (\cv{max\_poly},
default 3) and the operator set by \cv{operators}.  Optional extensions
include:
\begin{itemize}[nosep]
\item \emph{Nonlinear diffusion}: $\Delta(u^p)$ for specified powers~$p$;
\item \emph{Gradient nonlinearity}: $|\nabla u|^2 = u_x^2+u_y^2$;
\item \emph{Conservative forms}: $\partial_x(u^p),\;\partial_y(u^p)$.
\end{itemize}
The library is assembled by the fluent \cv{LibraryBuilder} API
in \cv{src/wsindy/library.py}.

\subsection{Multi-field library for coupled Vicsek--Morse dynamics}
\label{subsec:multifield_library}

When multiple coarse-grained fields are available
(see \S\ref{sec:multifield_cg} for field definitions), the pipeline
discovers a coupled PDE system.  For the Vicsek--Morse model with
density~$\rho$, polarisation flux $(p_x,p_y)$, and Morse
potential~$\Phi=W*\rho$, three separate regressions share a common
test-function bundle but have independent candidate libraries.

Table~\ref{tab:multifield_terms} lists the default term set for each
equation.  The Morse potential is precomputed via FFT convolution of the
density with the radial kernel
$W(r) = -C_r e^{-r/\ell_r} + C_a e^{-r/\ell_a}$
(see~\eqref{eq:morse_fr}).

\begin{table}[ht]
\centering
\caption{Default multi-field library terms for the Vicsek--Morse system
(\cv{src/wsindy/multifield.py}).}
\label{tab:multifield_terms}
\renewcommand{\arraystretch}{1.15}
\small
\begin{tabular}{lll}
\toprule
\textbf{Equation} & \textbf{Term name} & \textbf{Mathematical form} \\
\midrule
\multirow{6}{*}{$\rho_t = \cdots$}
  & \cv{div\_p}            & $\nabla\!\cdot\mathbf{p}$ \\
  & \cv{lap\_rho}          & $\Delta\rho$ \\
  & \cv{lap\_rho2}         & $\Delta(\rho^2)$ \\
  & \cv{lap\_rho3}         & $\Delta(\rho^3)$ \\
  & \cv{div\_rho\_gradPhi} & $\nabla\!\cdot(\rho\,\nabla\Phi)$ \\
  & \cv{lap\_p\_sq}        & $\Delta(|\mathbf{p}|^2)$ \\
\midrule
\multirow{8}{*}{$(p_x)_t = \cdots$}
  & \cv{px}                & $p_x$ \\
  & \cv{p\_sq\_px}         & $|\mathbf{p}|^2\,p_x$ \\
  & \cv{dx\_rho}           & $\partial_x\rho$ \\
  & \cv{dx\_rho2}          & $\partial_x(\rho^2)$ \\
  & \cv{lap\_px}           & $\Delta p_x$ \\
  & \cv{bilap\_px}         & $\Delta^2 p_x$ \\
  & \cv{rho\_dx\_Phi}      & $\rho\,\partial_x\Phi$ \\
  & \cv{p\_dot\_grad\_px}  & $\mathbf{p}\!\cdot\!\nabla p_x$ \\
\midrule
$(p_y)_t = \cdots$ & \multicolumn{2}{l}{\emph{(analogous to $p_x$ with $x\leftrightarrow y$)}} \\
\bottomrule
\end{tabular}
\end{table}

The $\rho$-equation terms encode diffusion, nonlinear diffusion, and
Morse-mediated drift; the momentum equations add linear drag ($p_x$),
cubic self-interaction ($|\mathbf{p}|^2 p_x$), pressure-like gradients
($\partial_x\rho$), viscosity ($\Delta p_x$), hyper-viscosity
($\Delta^2 p_x$), Morse forcing ($\rho\,\partial_x\Phi$), and
self-advection ($\mathbf{p}\!\cdot\!\nabla p_x$).

%
\section{Weak linear system assembly}
\label{sec:weak_system}

Given a density snapshot tensor $U\in\mathbb{R}^{T\times n_x\times n_y}$
and a collection of $K$~test functions $\{\psi_k\}_{k=1}^K$, the weak
formulation~\eqref{eq:ibp_temporal} produces the overdetermined linear
system
\begin{equation}
  \mathbf{b} = \mathbf{G}\,\mathbf{w},
  \qquad
  \mathbf{b}\in\mathbb{R}^K,\;
  \mathbf{G}\in\mathbb{R}^{K\times M},\;
  \mathbf{w}\in\mathbb{R}^M,
  \label{eq:weak_system}
\end{equation}
with components defined as follows.

\paragraph{Left-hand side (temporal derivative).}
Each entry of~$\mathbf{b}$ encodes the time derivative transferred
onto the test function:
\begin{equation}
  b_k = -\sum_{x,y,t}
    \frac{\partial\psi_k}{\partial t}(x,y,t)\;
    U(x,y,t)\;
    \Delta x\,\Delta y\,\Delta t.
  \label{eq:b_vector}
\end{equation}
The one-dimensional bump functions are defined analytically by
\eqref{eq:test_1d}, but the discrete derivative arrays entering
$\partial_t\psi_k$ are evaluated numerically on the support lattice via
one-dimensional finite-difference stencils and then assembled by
separable outer products, matching \cv{make\_separable\_psi}.

\paragraph{Right-hand side (library evaluations).}
For each library term $(\mathcal{L}_m,f_m)$, the $(k,m)$-entry of the
design matrix~$\mathbf{G}$ is
\begin{equation}
  G_{k,m}
  = \sum_{x,y,t}
    \bigl(\mathcal{L}_m[\psi_k]\bigr)(x,y,t)\;
    f_m\!\bigl(U(x,y,t)\bigr)\;
    \Delta x\,\Delta y\,\Delta t,
  \label{eq:G_matrix}
\end{equation}
where $\mathcal{L}_m[\psi_k]$ denotes the \emph{adjoint} of operator
$\mathcal{L}_m$ applied to the test function (e.g.\ $\Delta\psi_k$ for
the Laplacian, $-\partial_x\psi_k$ for $\partial_x$ under periodic BC).

\paragraph{Query-point selection.}
The $K$ test functions are centred at a uniform sub-grid with strides
$(s_t,s_x,s_y)$ (\cv{stride}, default $(2,2,2)$).  A temporal margin
of $\ell_t$~time steps is excluded at the beginning and end of the
time series to keep each test function's support within the data domain.
The resulting system is tall and thin ($K\gg M$), well-suited for
least-squares regression.

\paragraph{Implementation.}
All convolutions in \eqref{eq:b_vector}--\eqref{eq:G_matrix} are
evaluated via \cv{fft\_convolve3d\_same} (see \S\ref{subsec:fft_conv})
and then sampled at the query locations.  The full assembly is performed by
\cv{build\_weak\_system} in \cv{src/wsindy/system.py}.

\paragraph{Multiple trajectories.}
In production, WSINDy is fit from many training simulations sharing the same
parameter regime but differing in initial condition.  These trajectories are
\emph{not} concatenated end-to-start into one artificial long time series.
Instead, for each trajectory $c=1,\dots,C$ we assemble a separate weak system
\[
  \mathbf{b}^{(c)} = \mathbf{G}^{(c)}\mathbf{w},
\]
and then stack the resulting rows into one global regression problem,
\begin{equation}
  \begin{bmatrix}
    \mathbf{b}^{(1)}\\
    \mathbf{b}^{(2)}\\
    \vdots\\
    \mathbf{b}^{(C)}
  \end{bmatrix}
  =
  \begin{bmatrix}
    \mathbf{G}^{(1)}\\
    \mathbf{G}^{(2)}\\
    \vdots\\
    \mathbf{G}^{(C)}
  \end{bmatrix}
  \mathbf{w}.
  \label{eq:stacked_weak_system}
\end{equation}
This yields one common sparse coefficient vector for the scalar case.  In the
multi-field Vicsek--Morse setting, the same block-stacking idea is applied
separately to the $\rho$, $p_x$, and $p_y$ equations, producing three coupled
field-level regressions rather than a POD latent time-stepper.

%
\ifdefined\feedbackdraft
\section{Sparse regression and model selection}
\label{sec:regression}

\section{Uncertainty quantification}
\label{sec:uq}

\section{PDE forecasting}
\label{sec:forecast}
\else
\section{Sparse regression and model selection}
\label{sec:regression}

\subsection{Column preconditioning}
\label{subsec:precondition}

Before thresholding, each column of $\mathbf{G}$ is normalised to unit
$\ell^2$ norm:
\begin{equation}
  \tilde{G}_{:,j} = G_{:,j}\,/\,s_j,
  \qquad
  s_j = \lVert G_{:,j}\rVert_2,
  \label{eq:precondition}
\end{equation}
so that the thresholding criterion treats all candidate terms on equal
footing regardless of their physical scale.  After regression on the
normalised system $\tilde{\mathbf{G}}\tilde{\mathbf{w}}\approx\mathbf{b}$,
the physical coefficients are recovered by $w_j = \tilde{w}_j/s_j$.

\subsection{MSTLS: dominant-balance thresholding}
\label{subsec:mstls}

We use the \emph{Multi-Step Thresholded Least Squares} (MSTLS) algorithm
with the dominant-balance criterion of Minor et al.~\cite{minor2025mstls}.
Starting from all terms active, the algorithm iterates:
\begin{enumerate}[nosep]
\item Solve least squares on the active columns:
      $\tilde{\mathbf{w}}_{\mathcal{A}}
       = \arg\min_{\mathbf{w}}
         \lVert\tilde{\mathbf{G}}_{:,\mathcal{A}}\,\mathbf{w}
               -\mathbf{b}\rVert_2$.
\item Keep term~$j$ if its relative contribution lies in the band
      \begin{equation}
        \hat\lambda
        \;\le\;
        \frac{|\tilde{w}_j|\;\lVert\tilde{G}_{:,j}\rVert_2}
             {\lVert\mathbf{b}\rVert_2}
        \;\le\;
        \frac{1}{\hat\lambda}.
        \label{eq:dominant_balance}
      \end{equation}
\item Repeat until the active set~$\mathcal{A}$ converges
      (or a maximum of 25 iterations).
\end{enumerate}
The procedure is run for each threshold in a logarithmic grid
$\hat\lambda\in[10^{-4},10^1]$ (40 values), and the threshold
producing the lowest normalised loss is selected (ties broken
by sparsity).  Full pseudocode is given in
Algorithm~\ref{alg:mstls} (Appendix~\ref{app:mstls}).

\subsection{Automated model selection across test-function scales}
\label{subsec:model_selection}

The accuracy of the weak formulation depends critically on the test-function
half-widths $\boldsymbol\ell=(\ell_t,\ell_x,\ell_y)$: too narrow yields
noisy convolutions; too wide smooths out spatial structure.  Rather than
selecting $\boldsymbol\ell$ by hand, we sweep a log-spaced grid of
configurations (typically 5--20 trials) and score each with a composite
objective:
\begin{equation}
  S(\boldsymbol\ell)
  = L_{\mathrm{norm}}(\boldsymbol\ell)
  + \alpha\,\frac{n_{\mathrm{active}}(\boldsymbol\ell)}{M}
  + \beta\,\max\!\Bigl(0,\;
      \log_{10}\kappa(\boldsymbol\ell)
      - \log_{10}\kappa_0\Bigr),
  \label{eq:composite_score}
\end{equation}
where $L_{\mathrm{norm}}$ is the MSTLS normalised residual,
$n_{\mathrm{active}}$ the number of surviving terms, $M$ the total library
size, and $\kappa$ the condition number of the active sub-matrix.  The
default weights are $\alpha=0.1$ (complexity penalty), $\beta=0.01$
(conditioning penalty), and $\kappa_0=10^8$ (threshold above which
conditioning is penalised).

The configuration minimising~\eqref{eq:composite_score} is returned as the
best model.  A Pareto frontier in the two-dimensional space
$(L_{\mathrm{norm}},\;n_{\mathrm{active}})$ is also computed for
user inspection.

The model selection sweep, composite scoring, and Pareto-frontier
extraction are implemented in \cv{src/wsindy/select.py}
(\cv{wsindy\_model\_selection}).

%
\section{Uncertainty quantification}
\label{sec:uq}

\subsection{Bootstrap resampling}
\label{subsec:bootstrap}

To quantify coefficient uncertainty, we resample the $K$~rows of the weak
system $(\mathbf{b},\mathbf{G})$ uniformly at random with replacement,
refit MSTLS on each bootstrap replicate, and collect the resulting
coefficient vectors~\cite{messenger2021wsindy}:

\begin{enumerate}[nosep]
\item For $b=1,\dots,B$ (default $B=100$):
  \begin{enumerate}[nosep]
  \item Draw $K'=\lceil K\cdot f_{\mathrm{sub}}\rceil$ rows with
        replacement (default $f_{\mathrm{sub}}=1$).
  \item Refit MSTLS $\to$ $\mathbf{w}^{(b)}$, active set $\mathcal{A}^{(b)}$.
  \end{enumerate}
\item Aggregate:
  \begin{itemize}[nosep]
  \item Mean $\bar{w}_j = \frac{1}{B}\sum_b w_j^{(b)}$,\;
        Std $\sigma_j = \sqrt{\frac{1}{B}\sum_b(w_j^{(b)}-\bar{w}_j)^2}$.
  \item Inclusion probability
        $P_j = \frac{1}{B}\sum_b\mathbf{1}_{j\in\mathcal{A}^{(b)}}$.
  \item Percentile-based confidence intervals (default 95\%).
  \end{itemize}
\end{enumerate}
The implementation is in \cv{src/wsindy/uncertainty.py}
(\cv{bootstrap\_wsindy}).

\subsection{Stability selection}
\label{subsec:stability}

Stability selection~\cite{meinshausen2010stability} identifies which
discovered terms are robust to changes in the test-function scale and/or
the data subsample.  For each configuration~$c$ in a set of $n_c$~trials
(varying $\boldsymbol\ell$ and, optionally, bootstrap subsamples):
\begin{enumerate}[nosep]
\item Build the weak system for configuration~$c$.
\item Fit MSTLS and record the active set $\mathcal{A}^{(c)}$.
\end{enumerate}
The selection frequency of term~$j$ is
\begin{equation}
  f_j = \frac{1}{n_c}\sum_{c=1}^{n_c}\mathbf{1}_{j\in\mathcal{A}^{(c)}}.
  \label{eq:stability_freq}
\end{equation}
Terms are classified as
\emph{robust} ($f_j\ge 0.6$),
\emph{fragile} ($0<f_j<0.6$), or
\emph{never selected} ($f_j=0$).
The threshold 0.6 follows the recommendation of
Meinshausen \& B\"uhlmann~\cite{meinshausen2010stability}.
The implementation is in \cv{src/wsindy/stability.py}
(\cv{stability\_selection}).

%
\section{PDE forecasting}
\label{sec:forecast}

Once a sparse PDE has been identified, we integrate it forward in time
from a given initial condition to produce a density forecast.

\subsection{Strong-form RHS evaluation}
\label{subsec:rhs}

The right-hand side of the discovered PDE is evaluated pointwise in
physical space:
\begin{equation}
  u_t(x,y)
  = \sum_{m\in\mathcal{A}} w_m\;\mathcal{L}_m\!\bigl[f_m(u)\bigr](x,y),
  \label{eq:rhs_eval}
\end{equation}
where spatial derivatives are computed spectrally via FFT on the periodic
domain:
\[
  \widehat{\partial_x u}(k_x,k_y) = i\,k_x\,\hat{u}(k_x,k_y),
  \qquad
  \widehat{\Delta u}(k_x,k_y) = -(k_x^2+k_y^2)\,\hat{u}(k_x,k_y).
\]
This is implemented in \cv{src/wsindy/rhs.py} (\cv{wsindy\_rhs}).

\subsection{Linear--nonlinear splitting}
\label{subsec:splitting}

To enable exponential time-stepping, the discovered RHS is decomposed into
\begin{equation}
  u_t = L\,u + N(u),
  \label{eq:lin_nonlin_split}
\end{equation}
where $L$ collects all terms that are \emph{linear} in~$u$ (feature $=u$,
any operator---e.g.\ $\alpha\Delta u$, $\beta\partial_x u$) and $N(u)$
contains the nonlinear remainder (e.g.\ $\Delta u^2$, $|\nabla u|^2$).
In Fourier space, $L$ has a diagonal multiplier
\begin{equation}
  \hat{L}(k_x,k_y)
  = \sum_{m\in\mathcal{A}_{\mathrm{lin}}} w_m\,\sigma_m(k_x,k_y),
  \label{eq:L_hat}
\end{equation}
where $\sigma_m$ is the Fourier symbol of operator~$\mathcal{L}_m$
(e.g.\ $\sigma_{\Delta}=-(k_x^2+k_y^2)$, $\sigma_{\partial_x}=ik_x$).
The splitting is performed by \cv{split\_linear\_nonlinear} in
\cv{src/wsindy/integrators.py}.

\subsection{ETDRK4 integration}
\label{subsec:etdrk4}

When a linear component exists, we use the fourth-order exponential
time-differencing Runge--Kutta scheme (ETDRK4) of Cox \&
Matthews~\cite{cox2002etdrk}, stabilised with the contour-integral trick
of Kassam \& Trefethen~\cite{kassam2005etdrk4}.  The scheme treats the
linear part \emph{exactly} in Fourier space via matrix exponentials, so
that the time-step is not constrained by diffusive CFL limits.  The
nonlinear terms are integrated with fourth-order accuracy through four
sub-stage evaluations.  Full coefficient formulas and the per-step
algorithm are given in Appendix~\ref{app:etdrk4}.

When no linear terms are present, a classical fourth-order Runge--Kutta
(RK4) integrator is used as a fallback.

\paragraph{High-level interface.}
The function \cv{wsindy\_forecast} in \cv{src/wsindy/forecast.py}
selects the integrator automatically (\cv{method='auto'}): ETDRK4 when
linear terms exist, RK4 otherwise.  An optional density floor
(\cv{clip\_negative=True}) enforces $u\ge 0$ after each time step.
\fi
 
%
\begin{thebibliography}{10}

\bibitem{alvarez2025eqfree}
Hector~Vargas Alvarez, Dimitrios~G. Patsatzis, Lucia Russo, Ioannis Kevrekidis,
  and Constantinos Siettos.
\newblock Next generation equation-free multiscale modelling of crowd dynamics
  via machine learning, 2025.
\newblock arXiv:2508.03926.

\bibitem{bhaskar2019topology}
Dhananjay Bhaskar, Angelika Manhart, Jesse Milzman, John~T. Nardini, Kathleen
  Storey, Chad~M. Topaz, and Lori Ziegelmeier.
\newblock Analyzing collective motion with machine learning and topology.
\newblock {\em Chaos}, 29(12):123125, 2019.

\bibitem{black2020projection}
Felix Black, Philipp Schulze, and Benjamin Unger.
\newblock Projection-based model reduction with dynamically transformed modes,
  2020.
\newblock Published in ESAIM: Mathematical Modelling and Numerical Analysis,
  54(6):2011--2043, 2020.

\bibitem{siminos2015periodic}
Nazmi~Burak Budanur and Predrag Cvitanovi{\'c}.
\newblock Periodic orbit analysis of a system with continuous symmetry---a
  tutorial, 2015.
\newblock Published in Chaos, 25(7):073112, 2015.

\bibitem{budanur2015slices}
Nazmi~Burak Budanur, Predrag Cvitanovi{\'c}, Ruslan~L. Davidchack, and
  Evangelos Siminos.
\newblock Reduction of continuous symmetries of chaotic flows by the method of
  slices, 2015.
\newblock Published in Communications in Nonlinear Science and Numerical
  Simulation, 29(1-3):30--47, 2015.

\bibitem{dorsogna2006morse}
Maria~R. D'Orsogna, Yao-Li Chuang, Andrea~L. Bertozzi, and Lincoln~S. Chayes.
\newblock Self-propelled particles with soft-core interactions: Patterns,
  stability, and collapse.
\newblock {\em Physical Review Letters}, 96(10):104302, 2006.

\bibitem{krah2024robust_spod}
Philipp Krah, Arthur Baumg{\"a}rtner, Beata {\.Z}{\'o}rawski, and Julius Reiss.
\newblock A robust shifted proper orthogonal decomposition: Insights from
  proximal methods, 2024.

\bibitem{messenger2021wsindy}
Daniel~A. Messenger and David~M. Bortz.
\newblock Weak {SINDy} for partial differential equations.
\newblock {\em Journal of Computational Physics}, 443:110525, 2021.

\bibitem{peherstorfer2022kolmogorov}
Benjamin Peherstorfer.
\newblock Breaking the {Kolmogorov} barrier with nonlinear model reduction.
\newblock {\em Notices of the American Mathematical Society}, 69(5):725--733,
  2022.

\bibitem{reiss2018shifted_pod}
Julius Reiss, Philipp Schulze, J{\"o}rn Sesterhenn, and Volker Mehrmann.
\newblock The shifted proper orthogonal decomposition: A mode decomposition for
  multiple transport phenomena.
\newblock {\em SIAM Journal on Scientific Computing}, 40(3):A1322--A1344, 2018.
\newblock arXiv:1512.01985.

\bibitem{sonday2011alignment}
Benjamin~E. Sonday, Amit Singer, and Ioannis~G. Kevrekidis.
\newblock Noisy dynamic simulations in the presence of symmetry: Data alignment
  and model reduction, 2013.
\newblock Published in Annals of Applied Statistics, 7(3):1535--1557, 2013.

\bibitem{vicsek1995vicsek}
T.~Vicsek, A.~Czir{\'o}k, E.~Ben-Jacob, I.~Cohen, and O.~Shochet.
\newblock Novel type of phase transition in a system of self-driven particles.
\newblock {\em Physical Review Letters}, 75(6):1226--1229, 1995.

\end{thebibliography}

\end{document}
